\documentclass[oneside,a4paper,german,parskip=half,draft]{scrbook}
\usepackage[latin1]{inputenc}
\pdfcompresslevel=9

%Neue Deutsche Rechtschreibung
\usepackage{ngerman}

%Seitenheader
\usepackage{fancyhdr}
\pagestyle{fancy}
\rhead{\nouppercase{\rightmark}}
\lhead{\nouppercase{\leftmark}}

%Standartfont
\usepackage[T1]{fontenc}
%\usepackage{fourier}
\usepackage{libertine}

%Zusatzpaket f�r mathematische Ausdr�cke
\usepackage{amsmath}

%Zusatzfonts f�r mathbb usw.
\usepackage{amsfonts}
\usepackage{mathrsfs}

%Verbessertes Ref
\usepackage[german]{varioref}

%Links
\usepackage{hyperref}

%Stichwortverzeichnis
\usepackage{makeidx}
\makeindex

%Acronyme
\usepackage[nolist,nohyperlinks]{acronym}

%Anpassbare Enumerates/Itemizes
\usepackage{enumitem}

%Tikz/PGF Zeichnenpaket
\usepackage{pgf}
\usepackage{tikz}
\usetikzlibrary{mindmap,trees,decorations,decorations.pathreplacing,decorations.pathmorphing,calc,arrows,automata,through,backgrounds,shapes.geometric}

%F�r graphische Todos
\usepackage[colorinlistoftodos, obeyDraft]{todonotes}

%Paket zum Berechnen von Textbreiten und H�hen
\usepackage{calc}

%BibTeX
%\usepackage{cite}
%\usepackage{bibgerm}
%\bibliographystyle{gerplain}

%F�r das manipulieren von Captions in Figuren
\usepackage{caption}

%F�r das anpassen von theoremen
\usepackage{amsthm}

%F�r die farbigen Boxen
\usepackage{xcolor}
\usepackage{framed}

%Quellcode Einbindung
\usepackage{listings}

%Um die Seite in mehrere Spalten aufzuteilen
\usepackage{multicol}

%F�r Widerspruchszeichen (Blitz)
\usepackage{ulsy}

%F�r zwischenr�ume in W�rtern
\usepackage{xspace}

%F�r das umstellen zwischen Doppelten, EinEinhalben und Einzelnen Zeilenabst�nden
\usepackage{setspace}

%F�r Commandos mit zwei opotionalen Argumenten
\usepackage{twoopt}

{\catcode`\!=8 % funny catcode so ! will be a delimiter 
\catcode`\Q=3 % funny catcode so Q will be a delimiter 
\long\gdef\given#1{88\fi\Ifbl@nk#1QQQ\empty!} 
\long\gdef\blank#1{88\fi\Ifbl@nk#1QQ..!}% if null or spaces 
\long\gdef\nil#1{\IfN@Ught#1* {#1}!}% if null 
\long\gdef\IfN@Ught#1 #2!{\blank{#2}} 
\long\gdef\Ifbl@nk#1#2Q#3!{\ifx#3}% same as above 
}

%<Commandos>
%In geschweifte Klammern setzen
\newcommand{\gklamm}[1]{\ensuremath{\left\{#1\right\}}}

%In eckige Klammern setzen
\newcommand{\eklamm}[1]{\ensuremath{\left[#1\right]}}

%In runde Klammern setzen
\newcommand{\rklamm}[1]{\ensuremath{\left(#1\right)}}

%In Betragsstriche Setzen
\newcommand{\betrag}[1]{\ensuremath{\left|#1\right|}}

%Script zum Eingeben von l�ngeren Beispielen
\newcommand{\bsp}[3][]
{
\if\blank{#3}\textalign{\underline{#2}}{#1}\else
\underline{#2} #1
\par
\begingroup
\leftskip=1.28em
\setlist[1]{labelindent=1.28em, leftmargin=*}
#3
\par
\endgroup\fi
}

%Script um Text einzur�cken
\newcommand{\textalign}[2]{
\begin{minipage}[b]{\widthof{#1} + \widthof{\space}}
#1
\end{minipage}
\begin{minipage}[t]{\linewidth-\widthof{#1}-\widthof{\space}}
#2
\end{minipage}
}

%Script um Text einzur�cken ohne Ausgabe
\newcommand{\textfakealign}[3]{
\begin{minipage}[b]{\widthof{#1} + \widthof{\space}}
\if\blank{#2}$ $\else#2\fi
\end{minipage}
\begin{minipage}[t]{\linewidth-\widthof{#1}-\widthof{\space}}
#3
\end{minipage}
}

%Indexe
%Normaler Index
\newcommand{\indexn}[2][]{#2\if\blank{#1}\index{#2}\else\index{#1}\fi}
%Unterstrichener Index
\newcommand{\indexu}[2][]{\underline{#2}\if\blank{#1}\index{#2}\else\index{#1}\fi}
%Kursiver Index
\newcommand{\indexi}[2][]{\textit{#2}\if\blank{#1}\index{#2}\else\index{#1}\fi}
%Fetter Index
\newcommand{\indexb}[2][]{\textbf{#2}\if\blank{#1}\index{#2}\else\index{#1}\fi}

%Definitionen, Beispiele, S�tze
\newenvironment{fshaded}[2]{%
\def\FrameCommand{\fcolorbox{#1}{#2}}%
\MakeFramed {\FrameRestore}}%
{\endMakeFramed}

\theoremstyle{definition}
\newtheorem{definition}{\ac{Def.}}[chapter]
\newenvironment{fdefinition}[1][]{
\definecolor{shadecolor_definition}{rgb}{.85882,.89803,.94509}
%OLD: \definecolor{shadecolor_definition}{rgb}{.87,.94,.96}
\definecolor{framecolor_definition}{rgb}{.30980,.50588,.74117}
%OLD: \definecolor{framecolor_definition}{rgb}{0,0,0}
\begin{fshaded}{framecolor_definition}{shadecolor_definition}\begin{definition}[#1]}{\end{definition}\end{fshaded}}

\newtheorem{lemma}[definition]{Lemma}

\newtheorem{satz}{Satz}[chapter]
\newenvironment{fsatz}[1][]{
\definecolor{shadecolor_satz}{rgb}{.91764,.94509,.86666}
%OLD: \definecolor{shadecolor_satz}{rgb}{.87,.96,.56}
\definecolor{framecolor_satz}{rgb}{.60784,.73333,.34901}
%OLD: \definecolor{framecolor_satz}{rgb}{0,0,0}
\begin{fshaded}{framecolor_satz}{shadecolor_satz}\begin{satz}[#1]}{\end{satz}\end{fshaded}}

\newtheoremstyle{beweis}% name
	{3pt}			%Space above
	{3pt}			%Space below
	{}				%Body font
	{}				%Indent amount (empty = no indent, \parindent = para indent)
	{\sc}	%Thm head font
	{:}			%Punctuation after thm head
	{.5em}		%Space after thm head: " " = normal interword space; \newline = linebreak
	{}				%Thm head spec (can be left empty, meaning `normal')

\theoremstyle{beweis}
\newtheorem{beweis}{Beweis}[chapter]
\newenvironment{fbeweis}[1][]{
\definecolor{shadecolor_beweis}{rgb}{.94901,.85882,.85882}
%OLD: \definecolor{shadecolor_beweis}{rgb}{1.00,.96,.92}
\definecolor{framecolor_beweis}{rgb}{.75294,.31372,.30196}
%OLD: \definecolor{framecolor_beweis}{rgb}{0,0,0}
\begin{fshaded}{framecolor_beweis}{shadecolor_beweis}\begin{beweis}[#1]}{\end{beweis}\end{fshaded}}

\theoremstyle{remark}
\newtheorem{notation}{Notation}[chapter]

\newtheoremstyle{note}% name
	{3pt}			%Space above
	{3pt}			%Space below
	{}				%Body font
	{}				%Indent amount (empty = no indent, \parindent = para indent)
	{\bfseries}	%Thm head font
	{:}			%Punctuation after thm head
	{.5em}		%Space after thm head: " " = normal interword space; \newline = linebreak
	{}				%Thm head spec (can be left empty, meaning `normal')

\newtheoremstyle{notelite}% name
	{3pt}			%Space above
	{3pt}			%Space below
	{}				%Body font
	{}				%Indent amount (empty = no indent, \parindent = para indent)
	{\itshape}	%Thm head font
	{:}			%Punctuation after thm head
	{.5em}		%Space after thm head: " " = normal interword space; \newline = linebreak
	{}				%Thm head spec (can be left empty, meaning `normal')

\theoremstyle{note}
\newtheorem*{beispiel}{Beispiel}

\theoremstyle{notelite}
\newtheorem*{anmerkung}{Anmerkung}
\newtheorem*{beobachtung}{Beobachtung}
\newtheorem*{bekannt}{Bekannt}
\newtheorem*{idee}{Idee}
\newtheorem*{frage}{Frage}
\newtheorem*{problem}{Problem}
\newtheorem*{also}{Also}
\newtheorem*{bemerkung}{Bemerkung}
\newtheorem*{ziel}{Ziel}
%%%%%%%%Arabische in R�mische Zahl umwandeln
\newcommand{\RM}[1]{\ensuremath{\mbox{\MakeUppercase{\romannumeral #1}}}}
%</Commandos>

%<Abk�rzungen>
\newcommand{\Ra}{\ensuremath{\Rightarrow}}
\newcommand{\ra}{\ensuremath{\rightarrow}}
\newcommand{\La}{\ensuremath{\Leftarrow}}
\newcommand{\la}{\ensuremath{\leftarrow}}
\newcommand{\Lra}{\ensuremath{\Leftrightarrow}}
\newcommand{\lra}{\ensuremath{\leftrightarrow}}
\newcommand{\hra}{\ensuremath{\hookrightarrow}}
\newcommand{\mal}{\ensuremath{\cdot}}
\newcommand{\klgl}{\ensuremath{\leq}}
\newcommand{\grgl}{\ensuremath{\geq}}
\newcommand{\oder}{\ensuremath{\vee}}
\newcommand{\und}{\ensuremath{\wedge}}
\newcommand{\ggq}[1]{"`#1"'}
\newcommand{\tx}[1]{\ensuremath{\text{#1}}}
\newcommand{\Sigs}{\ensuremath{\Sigma^*}}
\newcommand{\SigB}{\ensuremath{\Sigma_{\tx{Bool}}}}
\newcommand{\SigBs}{\ensuremath{\SigB^*}}
\newcommand{\ub}[2]{\ensuremath{\underbrace{#1}_{#2}}}
\newcommand{\tc}[2][red]{\textcolor{#1}{#2}}
\newcommand{\R}{\ensuremath{\mb{R}}}
\newcommand{\N}{\ensuremath{\mb{N}}}
\newcommand{\seq}{\ensuremath{\subseteq}}
\newcommand{\fnewline}{\mbox{}\par}
\newcommand{\ural}[2]{\ensuremath{\displaystyle \matholongrightarrow_{#1}^{#2}}}
\newcommand{\rkl}[1]{\ensuremath{\rklamm{#1}}}
\newcommand{\un}{\ensuremath{\infty}}
\newcommand{\notsubseteq}{\ensuremath{\not\subseteq}}
\newcommand{\qacc}{\ensuremath{q_{\tx{accept}}}}
\newcommand{\qrej}{\ensuremath{q_{\tx{reject}}}}
\newcommand{\txcol}[2]{\textcolor{#1}{#2}}

%Underbrace mit viel text
\newcommand{\ubl}[3][]{
\ub{#2}{\begin{minipage}{\if\blank{#1}\widthof{\ensuremath{#2}}\else#1\fi}\begin{center}\scriptsize{#3}\end{center}\end{minipage}}
}

%Stack mit viel text
\newcommand{\stackl}[3][]{
\stack{\begin{minipage}{\if\blank{#1}\widthof{\ensuremath{#3}}\else#1\fi}\begin{center}\scriptsize{#2}\end{center}\end{minipage}}{#3}
}
%Sin, Cos, Tan, Dim, Span, Arccos
\DeclareMathOperator{\spano}{span}
\DeclareMathOperator{\bin}{Bin}
\DeclareMathOperator{\ld}{ld}
\DeclareMathOperator{\prog}{Prog}
\DeclareMathOperator*{\vd}{\vdash}
\DeclareMathOperator*{\matholongrightarrow}{\longrightarrow}
\DeclareMathOperator{\Konf}{Konf}
\DeclareMathOperator{\Def}{Def}
\DeclareMathOperator{\Code}{Code}

\newcommand{\sinx}[1]{\ensuremath{\sin{\left(#1\right)}}}
\newcommand{\cosx}[1]{\ensuremath{\cos{\left(#1\right)}}}
\newcommand{\tanx}[1]{\ensuremath{\tan{\left(#1\right)}}}
\newcommand{\dimx}[1]{\ensuremath{\dim{\left(#1\right)}}}
\newcommand{\spanx}[1]{\ensuremath{\spano{\left(#1\right)}}}
\newcommand{\arccosx}[1]{\ensuremath{\arccos{\left(#1\right)}}}
\newcommand{\stack}[2]{\ensuremath{\stackrel{#1}{#2}}}

%Todo, Insert, Mark, Img, Solved
\newcommand{\Todo}[1]{\todo[inline]{TODO: #1}}
\newcommand{\Insert}[1]{\todo[inline, color=green!40]{INSERT: #1}}
\newcommand{\Mark}[1]{\todo[inline, nolist, color=blue!40]{MARK: #1}}
\newcommand{\Img}[1]{\todo[inline, nolist, color=black!40]{IMG: #1}}
\newcommand{\Solved}[2][]{\todo[inline, nolist]{TODO: #2\if\empty{#1} -- DONE\else\\SOLVED: #1\fi}}

%Schriften
\newcommand{\mb}[1]{\ensuremath{\mathbb{#1}}}
\newcommand{\mc}[1]{\ensuremath{\mathcal{#1}}}

\newcommand{\matrixp}[1]{\ensuremath{\begin{pmatrix} #1 \end{pmatrix}}}

%Tikz Abk�rzungen
%%F�r Automaten Graphen
\newcommand{\tikzedge}[5]{\path[->] (#1) edge[#4] node[#5] {$#3$} (#2);}
\newcommand{\tikzloop}[4]{\tikzedge{#1}{}{#2}{loop #3}{#4};}

%%F�r Mathematische Graphen
%%%Coordinatensystem
\newcommandtwoopt{\tikzcoor}[4][x][y]{\draw[<->,thick] (0,#4) node[right] {$#2$} -- (0,0) -- (#3,0) node[below] {$#1$};}
%</Abk�rzungen>


%<Einstellung f�r Hyperref>%%%%%%%%%%%%%%%%%%%%%%%%%%%%%%%%%%%%%%%%%%%%%%%%%%%%
\hypersetup{colorlinks=false, linkcolor=black, breaklinks=true, bookmarksdepth=3,unicode=true,bookmarksnumbered=true,
pdftitle={Hefter f�r Grundlagen der Informatik-- Stand: \today},
pdfauthor={Thaller Alexander},
pdfsubject={Hefter f�r das Fach Grundlagen der Informatik aus der Vorlesung von Professor Dr. rer. nat. Jorg Striegnitz (Str) f�r das 2. Semester der Informatik im Sommersemester 2009 an der Hochschule Regensburg.},
pdfkeywords={informatik,studium,hefter,gi,sommersemester,2009,grundlagen der informatik}}
%</Einstellung f�r Hyperref>%%%%%%%%%%%%%%%%%%%%%%%%%%%%%%%%%%%%%%%%%%%%%%%%%%%

%<Dokument Daten>%%%%%%%%%%%%%%%%%%%%%%%%%%%%%%%%%%%%%%%%%%%%%%%%%%%%%%%%%%%%%%
\titlehead{\includegraphics[width=34mm]{FH-Logo.pdf}}
\subject{Informatik (Bachelor) 2. Semester}
\author{Alexander Thaller}
\title{Grundlagen der Informatik\footnote{Gefundene Fehler oder Verbesserungsvorschlge bitte hier \href{http://fhrein2ss09.codeplex.com}{http://fhrein2ss09.codeplex.com/} berichten oder alternativ mir eine E-Mail schreiben \href{mailto:alexander.thaller@stud.fh-regensburg.de}{alexander.thaller@stud.fh-regensburg.de}. Vielen Dank.}}
\subtitle{Hefter des Sommersemesters 2009}
\date{stand \today}
\publishers{Aus der Vorlesung von Professor Dr. rer. nat. Jorg Striegnitz (Str)}
%</Dokument Daten>%%%%%%%%%%%%%%%%%%%%%%%%%%%%%%%%%%%%%%%%%%%%%%%%%%%%%%%%%%%%%
\begin{document}
\frontmatter
\begin{acronym}
	\acro{acr}{Akronym}
	\acro{z.B.}[\ensuremath{\tx{z.\,B.}\xspace}]{zum Beispiel}
	\acro{Def.}{Definition}
	\acro{elem.}{elementige}
	\acro{Teilm.}{Teilmenge}
	\acro{Anz.}{Anzahl}
	\acro{z.z.}[\ensuremath{\tx{z.\,z.}\xspace}]{zu zeigen}
	\acro{bzw.}{beziehungsweise}
	\acro{etc.}{et cetera}
	\acro{d.h.}[\ensuremath{\tx{d.\,h.}\xspace}]{daher}
	\acro{OBdA}{Ohne Beschr�nkung der Allgemeinheit}
	\acro{differenz.}{differenzierbar}
\end{acronym}

%\maketitle
\setcounter{tocdepth}{1}
\setcounter{secnumdepth}{3}
\clearpage
\pdfbookmark[0]{\contentsname}{tocanc}
\tableofcontents
\newpage
\listoftodos
\newpage
%<Text>%%%%%%%%%%%%%%%%%%%%%%%%%%%%%%%%%%%%%%%%%%%%%%%%%%%%%%%%%%%%%%%%%%%%%%%%
%\section*{�bersicht}
\begin{description}
\item[Professor:] Jorg Striegnitz
\item[Homepage:] \href{http://homepages.fh-regensburg.de/~stj39817/index.html}{http://homepages.fh-regensburg.de/~stj39817/index.html}
\item[E-Mail:] \href{mailto:joerg.striegnitz@informatik.fh-regensburg.de}{joerg.striegnitz@informatik.fh-regensburg.de}
\item[Zeitplan:]\mbox{}\par
		\begin{itemize}
		\item 4 Stunden pro Woche
		\item 2 Stunden Vorlesung
		\item 2 Stunden �bungen
		\end{itemize}
\item[Klausur:] Zulassungsvorrausetzungen: 50\% in den �bungen und eine Aufgabe Vorrechnen
\item[Tutoren:]\mbox{}\par
		\begin{itemize}
		\item Herr Bauer (Mo. 4. Stunde)
		\item Herr Lainer (Do. 6. Stunde)
		\item Frau Sporrer (U4M)
		\item Frau Franziska Brandstetter (Di. 2. Stunde) (\href{mailto:franziska.brandstetter@web.de}{franziska.brandstetter@web.de})
		\item Herr Brillner (U612) (\href{mailto:billner.michael@web.de}{billner.michael@web.de})
		\end{itemize}
\end{description}

\mainmatter
\chapter{Grundlagen}
\Mark{Kapitel 2}
\begin{figure}[htb]
\centering
\begin{tikzpicture}[auto]
	\node[draw, rectangle, minimum width=3cm, minimum height=1.5cm] (start) at (1.5,0.75){};
	\node[below of=start, node distance = 1.25cm]{Tastatur};
	\foreach \x in {0.25, 1.00, 1.75, 2.50}
	{
		\foreach \y in{0.25, 1.00}
			\draw (\x,\y) rectangle (\x+0.25,\y+0.25);
	}
	\node[draw, rectangle, right of=start, minimum height=2.5cm, text width = 1.5cm, minimum width = 1.5cm, node distance = 3.5cm] (second) {10010101 00101000 01101101 01011101 10011010};
	\node[below of=second, node distance = 1.75cm]{Computer};
	\node[above of=second, node distance = 1.75cm](secondb){\LARGE{V}\normalsize{erarbeitung}};
	\node[left of=secondb, node distance = 3.5cm] {\LARGE{E}\normalsize{ingabe}};
	\node[right of=secondb, node distance = 3.5cm] {\LARGE{A}\normalsize{usgabe}};
	\node[draw, rectangle, right of=second, minimum height=2cm, text width = 2cm, minimum width = 2cm, node distance = 3.5cm] (third) {Lorem ipsum dolor sit amet\dots};
	\node[below of=third, node distance = 1.5cm]{Bildschirm};

	\draw[->] (start.east) -- (second.west);
	\draw[->] (second.east) -- (third.west);
\end{tikzpicture}
\end{figure}
\Img{GI(V)-17.03.2009-IMG1}

\begin{description}
\item[Ziel:] Formale Darstellung der Ein-/Ausgaben eines Computers
		\begin{itemize}[label=\Ra]
		\item Repr�sentation von Daten
		\end{itemize}
\end{description}

\begin{anmerkung}
Computerprogramme sind ebenfalls Daten.
\end{anmerkung}

\section{Alphabet}
\Mark{Section 2.1}
\begin{fdefinition}
Eine endliche Menge, die nicht leer ist, von \indexn[Symbol]{Symbolen} $\Sigma$ hei�t \indexn{Alphabet}\index{$\Sigma$ (Alphabet)}. Die Elemente von $\Sigma$ hei�en \indexn[Buchstabe]{Buchstaben} (\indexn{Zeichen}, Symbole). \qed
\Mark{Definition 2.1}
\end{fdefinition}

\begin{beispiel}
\mbox{}\par
\begin{itemize}
\item $\Sigma_{\text{RGB}}$ = \gklamm{\text{\underline{rot}}, \text{\underline{gr�n}}, \text{\underline{blau}}}
\item $\Sigma_{\text{Bool}}$ = \gklamm{0, 1} das Boole'sche Alphabet
\item $\Sigma_{\text{lat}}$ = \gklamm{a, b, c, \dots, z} das lateinische Alphabet
\item $\Sigma_{\text{Tastatur}}$ = \gklamm{A, B, \dots, Z, \text{\textvisiblespace}, <, >, (, ), \dots}
\item $\Sigma_{\text{m}}$ =  \gklamm{0, \dots, m -1} f�r $m \grgl 1$ die $m$-adische Darstellung einer Zahl.
\item $\Sigma_{\text{logic}}$ = \gklamm{0, 1, x, (, ), \oder, \und, \neg} ein Alphabet f�r logische Ausdr�cke
\end{itemize}
\end{beispiel}

\section{Wort}
\begin{fdefinition}[Wort]
Sei $\Sigma$ ein Alphabet. Ein \indexn[Wort]{Wort}\index{$\lambda$ (Wort)} �ber $\Sigma$ ist eine endliche (eventuell leere) Folge von Buchstaben. Das leere Wort $\epsilon$ (oder manchmal auch $\lambda$) ist die leere Buchstabenfolge. Die Menge aller W�rter �ber $\Sigma$ bezeichnen wir mit $\Sigma^*$. Ferner sei $\Sigma^+ = \Sigma^* - \gklamm{\epsilon}$. Fortan gehen wir davon aus, dass $\epsilon \notin \Sigma$ gilt. \qed
\Mark{Definition 2.2}
\end{fdefinition}

\begin{beispiel}
$\overbrace{(0, 1, 0, 1, 0, 1, 0)}^{\text{Z}}$ ist ein Wort �ber $\Sigma_{\text{Bool}}$, $\Sigma_{\text{Tastatur}}$ und $\Sigma_{\text{logic}}$.
\begin{itemize}
\item Z ist \underline{Wort} �ber $\Sigma_{\text{Bool}}$
\item Z ist aus $\Sigma_{Bool}^*$
\end{itemize}
$\epsilon$ ist ein Wort �ber einem \ac{bel.} Alphabet ($\epsilon \in \Sigma^*$ f�r alle Alphabete $\Sigma^*$)
\end{beispiel}

\begin{notation}
W�rter werden k�nftig ohne Kommata geschrieben und lassen auch die Klammern weg.\\
Also \ac{z.B.} 01010 statt (0,1,0,1,0)\\
Beachte: Unterschied Wort TI (Folge von Symbolen) und Wort \ggq{Deutsch}
\end{notation}

\subsection{Wortl�nge}
\begin{fdefinition}[Wortl�nge]
Sei $\Sigma$ ein Alphabet und $w \in \Sigma^*$. Die \indexn[Wort!Wortl�nge]{Wortl�nge} $\betrag{w}$ eines Wortes $w$ ist die L�nge des Wortes als Folge. F�r $a \in \Sigma$ bezeichnet $\betrag{w}_a$ die Anzahl der Vorkommen von $a$ in $w$. \qed
\Mark{Definition 2.3}
\end{fdefinition}

\begin{beispiel}
\mbox{}\par
\begin{itemize}
\item $\betrag{001001} = 6$
\item $\betrag{001001}_1 = 2$
\item $\betrag{\epsilon} = 0$
\item $\betrag{\text{\textvisiblespace}} = 1$
\end{itemize}
\end{beispiel}

\section{Konkatenation und Verkettung}
\begin{fdefinition}[Konkatenation/Verkettung]
Die \indexn[Alphabet!Konkatenation]{Konkatenation} f�r ein Alphabet $\Sigma$ ist eine Abbildung (oder auch Funktion) $K$:
\[\Sigma^* \times \Sigma^* \ra \Sigma^*\text{, so dass} K(u, v) = uv\]
f�r alle $u, v \in \Sigma^*$. Anstelle von $K(u, v)$ schreiben wir $u \mal v$. \qed
\Mark{Definition 2.4}
\end{fdefinition}

\begin{beobachtung}Konkatenation ist assoziativ. \Ra $(a \mal b) \mal c = a (b \mal c)$\\
Ferner gilt: $\epsilon \mal w = w = w \mal \epsilon$ Monoid.
\end{beobachtung}

\section{Induktive Definitionen}
\subsection{Induktive Definition des Wortbegriffes}
Die Menge aller W�rter �ber $\Sigma$ ($\Sigma^*$) ist induktiv definiert durch:
\begin{enumerate}
\item $\epsilon \in \Sigma^*$
\item sei $w \in \Sigma^*$ und $a \in \Sigma$, so $w \mal a \in \Sigma^*$\\
		\begin{tikzpicture}
			\node{}
			child[draw, rectangle]{
				child[draw, rectangle]{
					child[draw, rectangle]{
						child{node[draw, rectangle]{$\epsilon$}}
						child{node[draw, circle]{a}}}
					child{node[draw, circle]{b}}}
				child{node[draw, circle]{c}}};
		\end{tikzpicture}
\end{enumerate}

\subsection{Induktive Definition der Wortl�nge}
\begin{enumerate}
\item f�r $w = \epsilon$ sei $\betrag{w} = 0$
\item f�r $w = v \mal x$ ist $\betrag{w} = \betrag{v} + 1$
		\begin{align*}
			\betrag{\underbrace{\epsilon a b}_{v} \underbrace{c}_{x}} &= \betrag{\underbrace{\epsilon a}_{v} \underbrace{b}_{x}} + 1\\
			&= \betrag{\underbrace{\epsilon}_{v} \underbrace{a}_{x}} + 1 + 1\\
			&= \betrag{\epsilon} + 1 + 1 + 1\\
			&= 0 + 1 + 1 + 1\\
			&= 3
		\end{align*}
\end{enumerate}

\section{Kanonische Ordnung}
\index{Kanonische Ordnung}
Um alle W�rter aus $\Sigma^*$ systematisch aufzuz�hlen, ordnen wir Alphabet und W�rter.
\begin{fdefinition}[Kanonische Ordnung]
Sei $\Sigma = \gklamm{a_1, a_2, \dots, a_n}$ und $<: \Sigma \times \Sigma$ eine Ordnung auf $\Sigma$ mit $a_1 < a_2 < a_3 < \dots < a_n$. Wir definieren die kanonische Ordnung auf $\Sigma^*$ wie folgt:
\[u < v \text{ \ac{gdw.} } \betrag{u} < \betrag{v} \oder\]
\[\betrag{u} = \betrag{v} \und u = w a_i u_1 \und v w a_j u_2 \text{ f�r } u, v \in \Sigma^* w, u_1, u_2 \in \Sigma^* \text{ und } a_i < a_j (i < j)\] \qed
\Mark{Definition 2.5}
\end{fdefinition}

\section{Teilwort, Pr�fixe, Suffixe}
\begin{fdefinition}[Teilwort, Pr�fixe, Suffixe]
Sei $\Sigma$ ein Alphabet und seien $v, w \in \Sigma^*$
\begin{itemize}
\item $v$ hei�t \indexn[Wort!Teilwort]{Teilwort} von $w$ \ac{gdw.} $\exists s, t \in \Sigma^*$: $w = s v t$
\item $v$ hei�t \indexn[Wort!Suffix]{Suffix} von $w$ \ac{gdw.} $\exists s \in \Sigma^*$: $w = sv$\\
		\Img{GI(V)-25.03.2009-IMG1}
		\begin{tikzpicture}
			\draw (0,0) -- (3.5em,0);
			\draw (3.5em,0) -- (3.5em,1em) node[left, midway]{$v$};
			\draw (3.5em,1em) -- (0,1em) node[above, midway]{$w$};
			\draw (0,1em) -- (0,0) node[right, midway]{$s$};
			\draw (1.75em,0) -- (1.75em,1em);
		\end{tikzpicture}

\item $v$ hei�t \indexn[Wort!Pr�fix]{Pr�fix} von $w$ \ac{gdw.} $\exists t \in \Sigma^*$: $w = vt$\\
		\Img{GI(V)-25-03.2009-IMG2}
		\begin{tikzpicture}
			\draw (0,0) -- (3.5em,0);
			\draw (3.5em,0) -- (3.5em,1em) node[left, midway]{$t$};
			\draw (3.5em,1em) -- (0,1em);
			\draw (0,1em) -- (0,0) node[right, midway]{$v$};
			\draw (1.75em,0) -- (1.75em,1em);
		\end{tikzpicture}

\item $v \neq \epsilon$ hei�t \underline{echtes} Teilwort (Pr�fix/Suffix) von $w$ \ac{gdw.} $v \neq w$ und $v$ ist Teilwort (Pr�fix/Suffix) von $w$.
\end{itemize} \qed
\Mark{Definition 2.6}
\end{fdefinition}

\section{Wortumkehr}
\index{Wortumkehr}
\begin{fdefinition}[Wortumkehr]
Sei $\Sigma$ ein Alphabet und sei $w = a_1 \mal a_2 \mal a_3 \dots a_n$ mit $a_i \in \Sigma (1 \klgl i \klgl n)$ ein Wort. Dann ist die Wortumkehr $w^{\mathcal{R}}$ von $w$ \ac{def.} durch $w^{\mathcal{R}} = a_n a_{n - 1} \dots a_2 a_1$ \qed
\Mark{Definition 2.7}
\end{fdefinition}

\subsection{Alternativ: induktive Definition}
\begin{itemize}
\item $\epsilon^{\mathcal{R}} = \epsilon$
\item falls $v = w \mal a$, so $v^{\mathcal{R}} = a w^{\mathcal{R}}$
\end{itemize}

\begin{beispiel}
$(\overline{a}\overline{b}\overline{c})^{\mathcal{R}}$ $v = \overline{a}\overline{b}\overline{c}$ $w = \overline{a}\overline{b}$ $a = \overline{c}$
\begin{align*}
\Ra v^{\mathcal{R}} &= \overline{c} \mal \underbrace{(\overline{a}\overline{b})^{\mathcal{R}}} v = \overline{a}\overline{b} w = \overline{a} a = \overline{b}\\
\Ra v^{\mathcal{R}} &= \overline{c} \overline{b} \mal (\overline{a})^{\mathcal{R}} v = \epsilon \mal \overline{a} w = \epsilon a = \overline{a}\\
v^{\mathcal{R}} &= \overline{c} \mal \overline{b} \mal \overline{a} \mal (\epsilon)^{\mathcal{R}}\\
&= \overline{c} \mal \overline{b} \mal \overline{a} \mal \epsilon\\
&= \overline{c} \mal \overline{b} \mal \overline{a}
\end{align*}
\end{beispiel}

\section{Iteration}
\begin{fdefinition}[Iteration]
Sei $\Sigma$ ein Alphabet. F�r alle W�rter $w \in \Sigma^*$ und $i \in \mb{N}$ definieren wir die $i$-te Iteration $w^i$ als 
\[w^0 = \epsilon, w^1 = w\]
und
\[w^i = w \mal w^{i -1}\]
\Mark{Definition 2.8}
\end{fdefinition}

\begin{beispiel}
\begin{align*}
aabb \mal bbab &= aabbbbab\\
&= a^2 b^4 ab\\
ababc &= (ab)^2 c
\end{align*}
\end{beispiel}

\section{Sprachen}
\index{Sprache}
\index{Sprache!Konkatenation}
\index{Kleene-Stern}
\index{$L$ (Sprache)}
\begin{fdefinition}[Sprache, Konkatenation von Sprachen, Kleene-Stern]
Sei $\Sigma$ ein Alphabet. Eine \indexn{Sprache} $L$ ist eine Teilmenge von $\Sigma^*$ ($L \subseteq \Sigma^*$) (potentiell unendlich!). Das Komplement $L^{\mathcal{C}}$ der Sprache $L$ bez�glich $\Sigma^*$ ist die Sprache $\Sigma* - L$. $L_{\emptyset}$ ist die leere Sprache; $L_{\epsilon} = \gklamm{\epsilon}$. Seien $L_1$ und $L_2$ Sprachen, dann ist
\[L_1 \mal L_2 = L_1 L_2 = \gklamm{u \mal v \in \left(\Sigma_1 \cup \Sigma_2\right)^x \vert u \in L_1 \und v \in L_2}\]
die Konkatenation von $L_1$ und $L_2$. Ferner definieren wir f�r eine Sprache $L$
\[L^0 = L_{\epsilon} L^{i + 1} = L \mal L^i\]
und
\[L^* = \bigcup_{i \in \mb{N}} L^i \text{ und } L^+ = \bigcup_{i \in \mb{N} - \gklamm{0}} L^i = L \mal L^*\]
\Mark{Definition 2.9}
\end{fdefinition}

\begin{beispiel}
\mbox{}\par
\begin{multicols}{2}
\begin{itemize}
\item $L_1 = \emptyset$
\item $L_2 = \gklamm{\epsilon}$
\item $L_3 = \gklamm{\epsilon, aa, bb, ab, ba}$
\item $L_4 = \gklamm{a, b}^* = \gklamm{\epsilon, a, b, aa, ab, ba, bb, \dots}$
\item $L_5 = \gklamm{a}^* = \gklamm{\epsilon, a, aa, aaa, \dots} = \gklamm{a^0, a^1, a^2, \dots}$
\item $L_6 = \gklamm{a^p \vert p = \text{ Primzahl}}$
\item $L_7 = \gklamm{a^i b^{2i} a^i \vert i \in \mb{N}}$
\item $L_8 =$ alle syntaktisch korrekten \textit{JAVA}-Programme.
\item $L_9 =$ alle typ-korrekter \textit{JAVA}-Programme.
\item $L_{10} =$ alle terminierenden \textit{JAVA}-Programme.
\end{itemize}
\end{multicols}
\end{beispiel}

\begin{enumerate}
\item Sprachen geeignet zur Darstellung aller Ein/Ausgaben?
\item Reicht $\Sigma_{\text{Bool}}$?
\item Effizienter Test, ob $w \in L$ gilt. \fbox{\la Automaten}
\item Wie kann man Sprachen formal definieren? \fbox{\la Grammatiken}
\item Sind Sprachen aufz�hlbar? \fbox{\la Berechenbarkeit}\\
		Effizient m�glich. \fbox{\la Komplexit�tstheorie}
\end{enumerate}

\begin{lemma}
Seien $L_1$, $L_2$ und $L_3$ Sprachen �ber einem Alphabet $\Sigma$. Dann gilt
\[L_1 \mal L_2 \cup L_1 \mal L_3 = L_1 \mal (L_2 \cup L_3)\]
\end{lemma}

\begin{fbeweis}
\mbox{}\par
\begin{itemize}
\item $L_1 \mal L_2 \cup L_1 \mal L_3 \subseteq L_1 \mal (L_2 \cup L_3)$
		\begin{itemize}
		\item $L_1 \mal L_2 \subseteq L_1 \mal (L_2 \cup L_3)$
				\begin{align*}
				L_1 \mal L_2 &= \gklamm{uv \vert u \in L_1 \und v \in L_2}\\
				&\subseteq \gklamm{uv \vert u \in L_1 \und L_2 \cup L_3}\\
				&= L_1 \mal (L_2 \cup L_3)
				\end{align*}
		\item $L_1 \mal L_3 \subseteq L_1 \mal (L_2 \cup L_3)$\\
				\ra analog
		\end{itemize}

\item $L_1 \mal (L_2 \cup L_3) \subseteq L_1 \mal L_2 \cup L_1 \mal L_3$
		\begin{align*}
		&x \in L_1 \mal (L_2 \cup L_3) \text{ beliebig}\\
		\Lra & x \in \gklamm{uv \vert u \in L_1 \und v \in L_2 \cup L_3} \text{ (Definition Konkatenation)}\\
		\Lra & \exists u \in L_1 \und \exists v \in L_2 \cup L_3: x = uv \text{ (�bergang zur Pr�dikaten Logik)}\\
		\Lra & \exists u \in L_1 \und (\exists v \in L_2 \oder \exists v \in L_3): x = uv \text{ (�bergang Pr�dikaten Logik 2)}\\
		\Lra & (\exists u \in L_1 \und \exists v \in L_2: x = uv) \oder (\exists u \in L_1 \und \exists v \in L_3: x = uv)\\ &\text{(Distributivgesetz der Pr�dikaten Logik)}\\
		\Lra & x \in \underbrace{\gklamm{uv \vert u \in L_1 \und v \in L_2}}_{= L_1 \mal L_2} \oder x \in \underbrace{\gklamm{uv \vert u \in L_1 \und v \in L_3}}_{= L_1 \mal L_3}\\
		\Lra & x \in L_1 \mal L_2 \oder x \in L_1 \mal L_3\\
		\Lra & x \in L_1 \mal L_2 \cup L_1 \mal L_3
		\end{align*}
\end{itemize}
\end{fbeweis}

\section{Sprachen zur Beschreibung von Problemen}
\Mark{Section 2.3}
\subsection{Entscheidungsprobleme}
\Mark{Subsection 2.3.1}
\begin{fdefinition}
Das Entscheidungsproblem ($\Sigma$, $L$) f�r ein gegebenes Alphabet $\Sigma$ und eine Sprache $L \in \Sigma^*$ ist, f�r jedes $w \in \Sigma^*$ zu entscheiden, ob
\[w \in L \oder w \notin L\]
\[\Sigma_{\text{Input}}^* \longrightarrow \fbox{$\Sigma_{\text{Bool}}^* \ra \Sigma_{\text{Bool}}^*$} \longrightarrow \Sigma_{\text{Output}}^*\]
\end{fdefinition}

Ein Algorithmus l�st das Entscheidungsproblem ($\Sigma$, $L$), falls $\forall w \in \Sigma^*$
\[A(w) = \begin{cases}1 & w \in L\\0 & \text{sonst.}\end{cases}\]
Wir sagen auch, dass $A$ $L$ erkennt.

\begin{fdefinition}
Existiert f�r eine Sprache $L$ ein Algorithmus, der diese erkennt, so hei�t $L$ \underline{rekursiv}.
\end{fdefinition}

\begin{beispiel}
Ein wichtiges Entscheidungsproblem ist der Primzahltest.
\[\rklamm{\Sigma_{\text{Bool}}, \gklamm{w \in \Sigma_{\text{Bool}}^* \vert \text{Zahlwert }(w)\text{ ist Primzahl}}}\]
\end{beispiel}

\begin{beispiel}
Das Erf�llbarkeitsproblem.
\[\rklamm{\Sigma_{\text{logic}}m \gklamm{w \in \Sigma_{\text{logic}}^* \vert w \text{ erf�llbar}}}\]
\[\neg \rklamm{A_1 \oder A_2} \und A_3 \und A_4 \und \neg A_4\]
Ein ebenfalls wichtiges Entscheidungsproblem ist die Frage der �quivalenz von Algorithmen $A$ und $B$.
\[\rklamm{\Sigma_{\tx{Tastatur}}, \gklamm{B \vert B \tx{ ist �quivalent zu } A}}\]
$A$ ist �quivalent zu $B$, falls $A$ auf jede m�gliche Eingabe genauso reagiert wie $B$ und umgekehrt $A$.
\[\fbox{A} \stack{\text{Optimiert}}{\longrightarrow} \fbox{B}\]
\end{beispiel}

\subsection{\texorpdfstring{Aufz�hlungsalgorithmus f�r $L$}{Aufz�hlungsalgorithmus f�r L}}
\begin{fdefinition}
Sei $\Sigma$ ein Alphabet und $L \subseteq \Sigma^*$. $A$ ist ein Aufz�hlungsalgorithmus f�r $L$, falls $A$ f�r jede Eingabe $n \in \mb{N} - \rklamm{0}$ die kanonisch (siehe Kanonische Ordnung) ersten $n$ W�rter $w_0, w_1, \dots w_{n -1}$ der Sprache $L$ ausgibt.
\end{fdefinition}

\begin{lemma}
Sei $\Sigma$ ein Alphabet und $L \subseteq \Sigma^*$.
\begin{itemize}
\item $L$ ist rekursiv $\stackrel{\tx{\ac{gdw.}}}{\Lra}$ es existiert Aufz�hlungsalgorithmus f�r $L$
\item $L$ entscheidbar
\end{itemize}
\end{lemma}

\subsection{Funktionsberechnungen}
\begin{fdefinition}[Funktionsproblems]
Seien $\Sigma_D$ und $\Sigma_W$ zwei Alphabete. Wir sagen, dass ein Algorithmus $A$ eine Funktion $f: \Sigma_D^* \ra \Sigma_W^*$ berechnet (oder das Funktionsproblem $f$ l�st), falls
\[\forall w \in \Sigma_D^*: A(w) = f(w)\]
Existiert ein Algorithmus, der $f$ berechnet, so hei�t $f$ berechenbar.
\end{fdefinition}

\subsection{Relationsprobleme}
\begin{fdefinition}
Seien $\Sigma_1$ und $\Sigma_2$ Alphabete und $R \subseteq \Sigma_1^* \times \Sigma_2^*$ eine Relation. Ein Algorithmus $A$ berechnet $R$ (oder l�st das Relationsproblem $R$), falls f�r jedes $w \in \Sigma_1^*$ gilt:
\[\rklamm{w, A(w)} \in R\]
\end{fdefinition}

\section{Kolmogorov Komplexit�t}
\Mark{Section 2.4}
\begin{beispiel}[Wortl�nge vs. Informationsgehalt]
Das Wort
\[w = 01 01 01 01 01 01 01\]
ist zwar l�nger als das Wort
\[v = 10 10 11 00\]
kann aber durch $(01)^7$ kompakter dargestellt werden. Es reichen offenbar f�nf Zeichen um die Information in $w$ zu kodieren.
\end{beispiel}

Das Auffinden einer k�rzeren Darstellung f�r ein Wort nennt man \indexu{Komprimierung}.
\begin{description}
\item[Ziel:] Um den Informationsgehalt zweier W�rter zu vergleichen, k�nnten wir zun�chst ein Standard-Komprimierungs Verfahren suchen.\\
\item[Problem:] Wie findet man ein Verfahren, welches immer die kleinstm�gliche Darstellung garantiert?
\end{description}

\begin{beispiel}[Laufl�ngenkodierung]
Wir �berf�hren eine Bin�rzahl zun�chst in die Darstellung
\[b_1^{l_1} b_2^{l_2}\dots, b_i \in \Sigma_{\tx{Bool}}^*\]
wobei wir auch die Exponenten in Bin�rschreibweise darstellen. Von dieser Darstellung gehen wir wie folgt �ber zu einem Wort �ber $\Sigma_{\sharp} = \rklamm{0, 1, \sharp}$
\[l_1 \sharp b_1 \sharp l_2 \sharp b_2 \dots\]
Nach $\Sigma_{\tx{Bool}}$ kommen wir durch Anwendung eines Homomorphismus.\\
$\Sigma = \gklamm{0, 1, \sharp}$
\[\left.\begin{array}{l}
h (0) = 00\\
h (1) = 11\\
h (\sharp) = 10
\end{array} \right\} h: \Sigma_{\sharp} \ra \Sigma_{\tx{Bool}}^*\]
\[h (a_1 \dots a_n) = h(a_1) \mal h(a_2) \mal \dots \mal h(a_n)\]
\[k (00\sharp 1) = \begin{array}{ccccccc}
h(0) & \mal & h(0) & \mal & h(\sharp) & \mal & h(1)\\
00 && 00 && 01 && 11\\
0 && 0 && 10 && 1
\end{array}\]
$H: \Sigma_1^* \ra \Sigma_2^*$
\begin{itemize}
\item $h(\epsilon) = \epsilon$
\item $h(u \mal v) = h(u) \mal h(v)$
\end{itemize}
\end{beispiel}

\begin{beispiel}
Sei $w = (100)^{11} (101)^3 (01)^3 \betrag{w} = 43$\\
In unserer Zwischendarstellung ist dies
\[1011\sharp 100\sharp 11\sharp 101\sharp 11\sharp 01\sharp\]
und mit dem Homomorphismus $h(0) = 00$, $h(1) = 11$ und $h(\sharp) = 10$ gelangen wir zu
\[11 00\, 11 11\, 10 11\, 00 00\, 10 11\, 11 10 \dots \text{ mit } \betrag{w'} = 42\]
\end{beispiel}

F�r gro�e Exponenten kann es sich lohnen, die Laufl�ngenkodierung auch auf diese Anzuwenden. Zum Beispiel:
\[(011)^{1048576} = (011)^{2^{20}}\]

Diese Strategie kann man beliebig fortsetzen, um zum Beispiel W�rter der Bauarten $(01)^{2^{2^n}}$ oder $(01)^{2^{2^{2^n}}}$ zu kodieren.\\
Allgemeines kompressions Verfahren wird durch anderen Umstand ebenfalls schwer auffindbar.

\begin{beispiel}[Codierung durch Primfaktorzerlegung]
Jede Zahl k�nnen wir als Folge von Primfaktoren eindeutig darstellen: zum Beispiel $884736 = 2^1 \mal 3^3 \mal 4^7$. Unter Einbeziehung der Basis k�nnen wir dies darstellen als.
\[\bin(e_1) \sharp \bin(p_1) \sharp \bin(e_2) \sharp \dots\]
Verfahren der Primfaktorzerlegung und der Laufl�ngenkodierung sind \textbf{unvergleichbar}. Eines der beiden Verfahren ist Grunds�tzlich besser.
\end{beispiel}

\begin{fdefinition}
F�r alle W�rter $w \in \Sigma_{\tx{Bool}}^*$ ist die Kolmogorov-Komplexit�t $K_c(w)$ definiert als die bin�re L�nge \textit{des k�rzesten \lstinline{C++} Programmes}, welches $w$ generiert.
\end{fdefinition}

\begin{lemma}
Es existiert eine Konstante $d$, so dass f�r jedes $w \in \Sigma_{\tx{Bool}}^*$ gilt:
\[K_c (w) \klgl \betrag{w} + d\]
\end{lemma}

\begin{fbeweis}
F�r jedes $w$ aus $\Sigma_{\tx{Bool}}^*$ nehmen wir folgendes \lstinline{C++} Programm:
\[P(\underbrace{w}_{\tx{Zeichenkette}}) = \left\{
\begin{minipage}{82mm}
\begin{lstlisting}[language=C++]
int main(int args, char** args) {
	char * ws = w; //"abcdefg"
	count++ w;
}
\end{lstlisting}
\end{minipage}\right.\]
\end{fbeweis}

Regelm��ige W�rter haben eine geringere Kolmogorov-Komplexit�t. Betrachte \ac{z.B.} W�rter der Bauart $v_n = 0^n$ f�r $n \in \mb{N} \backslash \gklamm{0}$
\[P(n) = \left\{
\begin{minipage}{70mm}
\begin{lstlisting}[language=C++]
unsigned int n;
for(int i = 1; i <= n; i++)
	cout << "0";
\end{lstlisting}
\end{minipage}\right.\]
Bis auf die konstante $n$ sind alle $P(n)$ identisch. Zur Kodierung von $n$ als Bin�rzahl ben�tigt man $\lceil \ld (n + 1)\rceil$ Bits. Damit
\[K_c (v_n) \klgl \lceil \ld (n + 1) \rceil + c \klgl \lceil \ld n \rceil + c' = \lceil \ld \betrag{v_n} \rceil + c'\]

Betrachte $u_n = 0^{n^2}$ $n \in \mb{N} - \gklamm{0}$
\[P(n) \begin{cases}
\begin{minipage}{2.5cm}
\begin{center}
$\cancel{\mbox{\lstinline!m = n + n!}}$
\end{center}
\end{minipage}
\left|
\begin{minipage}{65mm}
\begin{lstlisting}
m = n;
m = m * m;
for(int i = 1; i <= m; i++)
	cout << "0";
\end{lstlisting}
\end{minipage}\right.
\end{cases}\]
\[K_c (u_n) \klgl \lceil \ld n \rceil + c \klgl \lceil \ld \sqrt{\betrag{u_n}} \rceil + c\]
\begin{align*}
&\betrag{u_n} = n^2\\
\Lra &n = \sqrt{\betrag{u_n}}
\end{align*}

\begin{fdefinition}
Sei $n \in \mb{N}$. Die Kolmogorov-Komplexit�t
\[K_c (n) = K_c(\bin(n))\tx{.}\]
\end{fdefinition}

\begin{lemma}
F�r jedes $n \in \mb{N} \backslash {0}$ existiert ein Wort $w_n \in \Sigma_{Bool}^n$, so dass
\[K_c (w_n) \grgl \betrag{w_n}\]
\textbf{Es gibt W�rter der L�nge $n$, die nicht komprimierbar sind.}
\end{lemma}

\begin{fbeweis}
Wir haben genau $2^n$ W�rter $x_1, \dots, x_{2^n} \in \Sigma_{\tx{Bool}}^n$. Sei f�r $i = 1, \dots, 2^n \prog(x_i)$ der Maschinencode des \lstinline{C++} Programmes von $x_i$ und sei $K_c (x_i) = \betrag{\prog(x_i)}$ die L�nge eines k�rzesten Algorithmus. Es ist klar, dass f�r $x_i \neq x_j$ gilt: 
\[\prog(x_i) \neq \prog(x_j)\]
wir haben also $2^n$ unterschiedliche Programme. Es gen�gt zu zeigen, dass eines dieser Programme die L�nge $n$ hat.

Jetzt hilft das kombinatorische Argument: es ist unm�glich $2^n$ verschiedene Maschinencodes der L�nge kleiner als $n$ zu haben, denn die Anzahl aller unterschiedlichen, nicht leeren, W�rter aus $\Sigma_{\tx{Bool}}^{n - 1}$ ist
\[\sum_{i = 1}^{n - 1} 2^i = 2^n - 2 < 2^n\]
Also muss es unter den $2^n$ Programmen mindestens eins der L�nge $n$ geben.
\end{fbeweis}

\begin{fsatz}
Seien $A$ und $B$ Programmiersprachen. Es existiert eine Konstante $c_{A, B}$, die nur von $A$ und $B$ abh�ngig ist, so dass $\forall w \in \Sigma_{\tx{Bool}}^*$
\[\betrag{K_{c_A} (w) - K_{C_B}} \klgl c_{A, B}\]
\end{fsatz}

\begin{fbeweis}
trivial
\end{fbeweis}

\begin{fdefinition}[Zuf�llig]
Ein Wort $w \in \Sigma_{\tx{Bool}}^*$ hei�t zuf�llig, falls $K_C (w) \grgl \betrag{w}$. Eine Zahl $n$ hei�t zuf�llig, falls
\[K_c(n) = K_c (\bin(n)) \grgl \lceil \ld(n + 1) \rceil - 1\]
\end{fdefinition}

%\chapter{Endliche Automaten}
\Mark{Kapitel 3}
\section{Deterministische endliche Automaten (DEA)}
\begin{fdefinition}[\ac{DEA}]
Ein \ac{DEA} ist ein Quintupel $M = (Q, \Sigma, \delta, q_0, F)$ aus
\begin{itemize}
\item einer endlichen \indexb[Endliche Automaten!$Q$ Menge von Zust�nden]{Menge von Zust�nden} $Q$
\item einem \indexb[Endliche Automaten!$\Sigma$ Eingabealphabet]{Eingabealphabet} $\Sigma$
\item einer \indexb[Endliche Automaten!$\delta$ Transitionsfunktion]{Transitionsfunktion} $\delta$: $Q \times \Sigma \ra Q$
\item einen \indexb[Endliche Automaten!$q_0$ Startzustand]{Startzustand} $q_0 \in Q$
\item eine \indexb[Endliche Automaten!$F$ Menge von akzeptierenden Zust�nden]{Menge von akzeptierenden Zust�nden} $F \subseteq Q$\\
		\begin{anmerkung}
		$F$ kann leeres Element sein also kann der Automat auch keine Zust�nde annehmen!
		\end{anmerkung}
\end{itemize}
\Mark{Definition 3.1}
\end{fdefinition}

\begin{notation}
Anstelle von $\delta(q, a) = q$ schreiben wir
\[p \stack{a}{\longrightarrow}_M q\]
\end{notation}

\begin{beispiel}
$M^*$
\begin{center}
\begin{tikzpicture}[->,>=stealth',shorten >=1pt,auto,node distance=2.8cm,semithick]
	\node[initial,accepting,state,initial text=] (q0) {$q_0$};
	\node[state] (q1) [right of=q0] {$q_1$};
	\node[state] (q2) [below of=q0] {$q_2$};
	\node[accepting,state] (q3) [right of=q2] {$q_3$};

	%q0
	\path[->]
	(q0) edge [bend right] node [below] {0} (q1)
	(q0) edge [bend left] node [right] {1} (q2);

	%q1
	\path[->]
	(q1) edge [bend right] node [above] {0} (q0)
	(q1) edge [bend left] node [right] {1} (q3);

	%p2
	\path[->]
	(q2) edge [bend left] node [above] {0} (q3)
	(q2) edge [bend left] node [left] {1} (q0);

	%q3
	\path[->]
	(q3) edge [bend left] node [below] {0} (q2)
	(q3) edge [bend left] node [left] {1} (q1);
\end{tikzpicture}
\end{center}
\Img{GI(V)-15.04.2009-DIA-1}

$M = \rklamm{\gklamm{q_0, q_1, q_2, q_3}, \Sigma_{\tx{Bool}}m, \delta, q_0, \gklamm{q_0, q_3}}$ mit $\delta$:
\begin{table}[htb]
\begin{tabular}{c|c|c}
& $0$ & $1$\\\hline
$q_0$ & $q_1$ & $q_2$\\
$q_1$ & $q_0$ & $q_3$\\
$q_2$ & $q_3$ & $q_0$\\
$q_3$ & $q_2$ & $q_1$
\end{tabular}
\end{table}
\end{beispiel}

\begin{fdefinition}
Sei $M = (Q, \Sigma, \delta, q_0, F)$ ein \ac{DEA}. Eine \indexb{Konfiguration} von $M$ ist ein Paar aus $Q \times \Sigma^*$. $(q_0, w)$ hei�t \indexb{Startkonfiguration} und $(q, \epsilon)$ \indexn{Endkonfiguration} (f�r \ac{bel.} $w \in \Sigma^*$). Gilt f�r eine Endkonfiguration $(q, \epsilon)$ das $q \in F$, so hei�t diese akzeptierend, ansonsten verwerfend. \qed
\Mark{Definition 3.2}
\end{fdefinition}

\begin{fdefinition}
Sei $M = (Q, \Sigma, \delta, q_0, F)$ ein \ac{DEA}. Ein \indexb{Konfigurations�bergang} (ein Schritt) ist eine Relation $\vd_M$: $(Q \times \Sigma^*) \times (Q \times \Sigma^*)$, die \ac{def.} ist durch
\[(p, \underbrace{a}_{\tx{Lesekopf}}\overbrace{w}^{\tx{Band}}) \vd (q, w) \tx{ \ac{gdw.} } \delta(p, a) = q\]

\begin{center}
\begin{tikzpicture}[->,>=stealth',shorten >=1pt,auto,node distance=2.8cm,semithick]
\node[state] (nodeone) {$q$};
\node[state] (nodetwo) [right of=nodeone] {$q$};

\path[->] (nodeone) edge node[above] {$a$} (nodetwo);
\end{tikzpicture}
\end{center}

mit $p, q \in Q$, $w \in \Sigma^*$ und $a \in \Sigma$. $\vd_u$ nennen wir \indexb{Schrittrelation}.
\Mark{Definition 3.3}
\end{fdefinition}

\begin{fdefinition}
Sei $M = (Q, \Sigma, \delta, q_0, F)$ ein \ac{DEA} und $w \in \Sigma^*$. Eine \indexb{Berechnung} $c$ von $M$ ist eine endliche Folge von Konfigurationen
\[c = c_0, \dots, c_n \tx{ mit } c_i \vd c_{i + 1} (1 \klgl i \klgl n - 1)\]
$c$ ist eine Berechnung von $M$ f�r eine Eingabe $w$, falls $c_0 = (q_0, w)$ (Startkonfiguration) und $c_n = (q, \epsilon)$ (Endkonfiguration). Falls $c_n$ \ac{akzept.} Endkonfiguration ist, so ist C \ac{akzept.} Berechnung; ansonsten verwerfende Berechnung.
\Mark{Definition 3.4}
\end{fdefinition}

\begin{beispiel}
\[(q_0, \underline{1}001) \vd (q_2, 001) \vd (q_3, 01) \vd (q_2, 1) \vd (\underbrace{q_0}_{\in F}, \epsilon)\]
ist \ac{akzept.} Berechnung von $M^*$ (zu GI(V)-15.04.2009-DIA-1) auf $1001$
\end{beispiel}

\begin{fdefinition}
Sei $M = (Q, \Sigma, \delta, q_0, F)$ ein \ac{DEA}. Wir definieren: $\vd_M^*$: $(Q \times \Sigma^*) \times (Q \times \Sigma^*)$ als reflexiven und transitiven Abschluss von $\vd$:
\begin{itemize}
\item $(q, w) \vd^* (q, w)$ $\forall q \in Q, w \in \Sigma^*$
\item $(p, uv) \vd^* (q, v)$, falls $u = a_1 \dots a_n (a_i \in \Sigma)$ $\exists q_1, \dots, q_{n - 1} \in Q$, so dass
		\[(p, a_1 \dots a_n v) \vd (q_1, a_1 \dots a_n v) \vd \dots \vd (q_{n - 1}, a_n v) \vd (q, v)\]
		f�r $p, q \in Q$ $u, v \in \Sigma^*$
\end{itemize}
Die Fortsetzung $\hat{\delta}$: $Q \times \Sigma^* \ra Q$ der Transitionsfunktion $\delta$ auf W�rter \ac{def.} wir induktiv durch
\begin{itemize}
\item $\hat{\delta} (q, \epsilon) = q$
\item $\hat{\delta} (q, wa) = \delta(\hat{\delta} (q, w))$ $\forall w \in \Sigma^*$, $a \in \Sigma$ und $q \in Q$
\end{itemize}
\Mark{Definition 3.5}
\end{fdefinition}

\begin{notation}
Falls $\hat{\delta} (q, w) = p$, so schreiben wir
\[q \stack{w}{\longrightarrow} p\]
und sagen es existiert ein Pfad von $q$ nach $p$ mit Beschriftung $w$.
\end{notation}

\begin{fdefinition}
Sei $M = (Q, \Sigma, \delta, q_0, F)$ ein \ac{DEA}. Die von $M$ akzeptierte Sprache $L(M)$ ergibt sich aus:
\begin{align*}
L(M) &= \gklamm{w \in \Sigma^* \vert \hat{\delta}(q_0, w) \in F}\\
&= \gklamm{w \in \Sigma^* \vert q_0 \stack{w}{\longrightarrow} q \tx{ und } q \in F}\\
&= \gklamm{w \in \Sigma^* \vert (q_0, w) \vd^* (q, \epsilon) \tx{ mit } q \in F}
\end{align*}
Sprachen, die von einem \ac{DEA} erkannt werden nennt man \indexb{regul�re Sprachen}. Zwei \ac{DEA}s $M_1$ und $M_2$ sind �quivalent \ac{gdw.} $L(M_1) = L(M_2)$
\Mark{Definition 3.6}
\end{fdefinition}

\begin{lemma}
Der \ac{DEA} $M_x = \rklamm{\Sigma_{\tx{Bool}}, Q, \delta, q_0, F}$ erkennt die Sprache
\[L(M_x) = \gklamm{w \in \Sigma_{\tx{Bool}}^* \vert \rklamm{\betrag{w}_0 + \betrag{w}_1} \mod 2 = 0}\]
\end{lemma}

\begin{fbeweis}
F�r jedes Wort $w \in \Sigma_{\tx{Bool}}^*$ bleibt der Automat $M_x$ in einem bestimmten Zustand stehen. Betrachten wir f�r einen \ac{bel.} Zustand die Menge der W�rter, f�r die $M_x$ in exakt diesem Zustand terminiert.
\[\mathcal{K} (q) = \gklamm{w \in \Sigma_{\tx{Bool}}^*\vert q_0 \stack{w}{\longrightarrow} q}\]
Da $M_x$ deterministisch ist, kann ein Wort, welches in $\mathcal{K}(q)$ liegt nicht gleichzeitig in $\mathcal{K}\rklamm{q'}$ liegen (mit $q \neq q'$) - $\mathcal{K}(q) \cap \mathcal{K}\rklamm{q'} = \emptyset$. Die Menge der von $M_x$ erkannten W�rter entspricht der Vereinigung aller Mengen $\mathcal{K}(q)$, die sich auf akzeptierende Zust�nde beziehen.
\[L(M_x) = \bigcup_{q \in F} \mathcal{K}(q)\]
$\mathcal{K}$ ist eine �quivalenzklasse folgender Relation $R_s$
\begin{align*}
&u R_{\delta} v = \hat{\delta} (q_0, u) = \hat{\delta} (q_0, u)\\
&R_{\delta} (w, v)
\end{align*}
in $\mathcal{K} (q)$ befindet sich nur solche W�rter, die \ac{bzgl.} $R_{\delta}$ in Relation stehen.

Beweisidee:
\begin{itemize}[label=-]
\item Ordne $q_i (0 \klgl i \klgl 3)$ konkrete �quivalenz-Klassen zu \textcircled{$\star$}
\item Zeige $\mathcal{K} (q_0) \cup \mathcal{K} (q_3) = L(u)$
\item Zeige, dass \textcircled{$\star$} korrekt.
\end{itemize}
\Solved{�ndere \textcircled{*} in ein sch�neres Symbol}{Replaced with \textcircled{$\star$}}
Wir stellen folgende Induktionsannahme auf:
\begin{align*}
&\mathcal{K} (q_0) = \gklamm{w \in \Sigma_{\tx{Bool}}^*~\vert~\betrag{w}_0 \tx{ und } \betrag{w}_1 \tx{ sind gerade}}\\
&\mathcal{K} (q_1) = \gklamm{w \in \Sigma_{\tx{Bool}}^*~\vert~\betrag{w}_0 \tx{ ungerade } \vert \betrag{w}_1 \tx{ gerade}}\\
&\mathcal{K} (q_2) = \gklamm{w \in \Sigma_{\tx{Bool}}^*~\vert~\betrag{w}_0 \tx{ gerade } \vert \betrag{w}_1 \tx{ ungerade}}\\
&\mathcal{K} (q_3) = \gklamm{w \in \Sigma_{\tx{Bool}}^*~\vert~\betrag{w}_0 \tx{ und } \betrag{w}_1 \tx{ ungerade}}
\end{align*}

\begin{enumerate}
\item Induktionsanfang: wir zeigen, dass $\vert A$ f�r W�rter der L�nge $\klgl 2$ gilt
		\begin{align*}
		0:~~&\delta \rklamm{q_0, \epsilon} = q_0 \Ra \epsilon \in \mathcal{K}(q_0) &\checkmark\\
		1:~~&\delta (q_0, 1) = q_2, 1 \in \mathcal{K}(q_2) &\checkmark\\
			&\delta (q_0, 0) = q_1, 0 \in \mathcal{K}(q_1) &\checkmark\\
		2:~~&\hat{\delta} (q_0, 00) = q_0, 00 \in \mathcal{K} (q_0) &\checkmark\\
			&\hat{\delta} (q_0, 01) = q_3, 01 \in \mathcal{K} (q_3) &\checkmark\\
			&\hat{\delta} (q_0, 10) = q_3, 10 \in \mathcal{K} (q_3) &\checkmark\\
			&\hat{\delta} (q_0, 11) = q_0, 11 \in \mathcal{K} (q_0) &\checkmark
		\end{align*}
\item Induktionsschluss: Wir setzen voraus, dass der Induktionsanfang f�r W�rter $\betrag{u} \klgl i$ gilt. Wir wollen jetzt zeigen, dass der Induktionsanfang dann auch f�r W�rter der L�nge $i + 1$ gilt.\\
		Sei also $w \in \Sigma_{\tx{Bool}}^{i + 1}$, also $w = u \mal a$ mit $u \in \Sigma_{\tx{Bool}}^{i}$ und $a \in \Sigma_{\tx{Bool}}$. Wir m�ssen vier F�lle unterscheiden:
		\begin{enumerate}[label=\arabic*.]
		\item $\betrag{u}_0$ gerade, $\betrag{u}_1$ gerade. Laut Induktionsvoraussetzung gilt $\hat{\delta} (q_0, u) = q_0$ \ac{bzw.} $u \in \mathcal{K} (q_0)$. Wir erhalten
				\[\hat{\delta} (q_0, ua) = \delta \rklamm{\hat{\delta} (q_0, u), a} \stack{\tx{IV}}{=} \delta(q_0, a) = \begin{cases}q_1 \tx{ falls } a = 0\\q_2 \tx{ falls } a = 1\end{cases}\]
				Da $\betrag{u1}_0$ gerade und $\betrag{u1}_1$ ungerade entspricht $\hat{\delta} (q_0, u1) = q_2$ der IA.\\
				Da $\betrag{u0}_0$ ungerade und $\betrag{u0}_1$ gerade entspricht $\hat{\delta} (q_0, u0) = q_1$ der IA.
		\item analog, trivial, selbst! 
		\item analog, trivial, selbst!
		\item analog, trivial, selbst!
		\end{enumerate}
\end{enumerate}
\end{fbeweis}

\section{Nichtdeterministische endliche Automaten (NEAs)}
\Mark{Section 3.3}
\begin{center}
\begin{tikzpicture}[->,>=stealth',shorten >=1pt,auto,node distance=2.8cm,semithick]
	\node[initial,state,initial text=] (A) {};
	\node[state] (B) [right of=A] {};
	\node[accepting,state] (C) [right of=B] {};

	\path (A) edge [loop above] node {$0,1$} ();
	\path (A) edge node {$0$} (B);
	\path (B) edge node {$0$} (C);
\end{tikzpicture}
\end{center}
\Img{IMG-GI-09-04-21-1}

\begin{fdefinition}
Ein \ac{NEA} $M$ ist ein Quintupel $(Q, \Sigma, \delta, q_0, F)$ aus
\begin{itemize}
\item $Q, \Sigma, q_0, F$ wie \ac{DEA}
\item eine Transitionsfunktion $\delta$: $Q \times \Sigma \ra \mathcal{P} (Q)$
\end{itemize}
\end{fdefinition}

\begin{notation}
Falls $q \in \delta(p, a)$ schreiben wir $p \stack{a}{\ra} q$
\end{notation}

\begin{fdefinition}
Konfiguration analog zu \ac{DEA} $Q \times \Sigma^*$
\end{fdefinition}

\begin{fdefinition}
Schrittrelation kann ebenfalls �bernommen werden, da $p \stack{a}{\ra} q$ f�r \ac{NEA} angepasst.
\[(p, aw) \vd (q, w) \tx{\ac{gdw.}} p \stack{a}{\ra} q\]
\end{fdefinition}

\begin{fdefinition}
Sei $M = \rklamm{Q, \Sigma, \delta, q_0, F}$ eine \ac{NEA}. Eine Berechnung $C$ von $M$ ist eine Folge von Konfigurationen
\[C = C_0, \dots, Q_n \tx{ mit } C_i \vd C_{i + 1} (0 \klgl i \klgl n - 1)\]
$C$ ist eine Berechnung von $M$ f�r die Eingabe $w$, falls $C_0 = (q_0, w)$ Startkonfiguration ist und entweder
\begin{itemize}
\item $C_n = (q, \epsilon)$ eine Endkonfiguration ist, oder
\item $C_n = (q, aw)$ und $\delta(q, a) = \emptyset$ gilt (der Automat \ggq{st�rzt ab})
\end{itemize}
\ac{akzept.} und verwerfende \ac{Ber.} analog \ac{DEA}. Reflexiver und transitiver Abschluss von $\vd$ analog \ac{DEA}.
\[(q,w) \vd^* (q,w)\]
\[(q, a_1, a_2, a_3) \vd^* (q', \epsilon)\]
\Mark{Definition 3.8}
\end{fdefinition}

\begin{fdefinition}
Sei $M = (Q, \Sigma, \delta, q_0, F)$ ein \ac{NEA}. Die Fortsetzung $\hat{\delta}: Q \times \Sigma^* \ra \mathcal{P}(Q)$ von $\delta$ auf W�rter ist \ac{def.} durch
\begin{itemize}
\item $\hat{\delta} (q, \epsilon) = \gklamm{q}$
\item Sei $w = ua$ und $\hat{\delta} (q, u) = \gklamm{q_1, \dots, q_n}$, dann gilt
		\[\hat{\delta} (q, ua) = \bigcup_{i = 1}^n \delta (q_i, a)\]
		\[w, u \in \Sigma^*, a \in \Sigma, q_i, q \in Q, 1 \klgl i \klgl n\]
\end{itemize}
\Mark{Definition 3.9}
\end{fdefinition}

\begin{fdefinition}
Sei $M = (Q, \Sigma, \delta, q_0, F)$ ein \ac{NEA}. Die von $M$ \ac{akzept.} Sprache $L(M)$ ergibt sich aus
\begin{align*}
L(M) &= \gklamm{w \in \Sigma^* \vert (q_0, w) \vd^{*} (q, \epsilon) \tx{ mit } q_f \in F}\\
&= \gklamm{w \in \Sigma^* \vert \hat{\delta} (q_0, w) \cap F \neq \emptyset}\\
&= \gklamm{w \in \Sigma^* \vert q_0 \stack{w}{\ra} q_f \tx{ und } q_f \in F}
\end{align*}

\begin{center}
\begin{tikzpicture}[->,>=stealth',shorten >=1pt,auto,node distance=2.8cm,semithick]
\node[state, initial, initial text=] (q0) {$q_0$};
\node[state] (q1) [right of=q0] {$q_1$};
\node[state, accepting] (q2) [right of=q1] {$q_2$};

\path (q0) edge [loop above] node[left]{$1,0$} ();
\path (q2) edge [loop above] node[right]{$1,0$} ();

\path (q0) edge node[above] {$1$} (q1);
\path (q1) edge node[above] {$1$} (q2);
\end{tikzpicture}
\end{center}
\Mark{Definition 3.10}
\end{fdefinition}

\begin{lemma}
$L(M_{11}) = \overbrace{\gklamm{w \in \Sigma_{\tx{Bool}}^* \vert w = u11v \tx{ mit } u,v \in \Sigma_{\tx{Bool}}^*}}^L$\\
Wie �blich erfolgt der \ac{Bew.} in 2 Schritten
\begin{itemize}
\item $\gklamm{u11v \vert u, v \in \Sigma_{\tx{Bool}}^*} \subseteq L (M_{11})$. Sei $w \in L$; \ac{d.h.} $w = u11v$ f�r \ac{bel.} $u, v \in \Sigma_{\tx{Bool}}^*$. Zu zeigen: es gibt \ac{akz.} Berechnung von $M_{11}$ f�r $w$. Da $q_0 \in \delta(q_0, 0) \cap \delta (q_0, 1)$ \ac{ex.} f�r jedes $z \in \Sigma_{Bool}^*$ eine Berechnung
		\[(q_0, z) \vd^* (q_0, \epsilon)\]
		gleicherma�en
		\[(q_2, z) \vd^* (q_2, \epsilon), \tx{ da } q_2 \in \delta (q_2, 0) \cap \delta(q_2, 1)\]
		Damit erhalten wir f�r \ac{Ber.} f�r $w$ wie folgt
		\[(q_0, u11v) \vd^* (q_0, 11v) \vd (q_1, 1v) \vd (q_2, v) \vd^* (q_2, \epsilon)\]

\item $L(M_{11}) \subseteq L.$ Sei $w \in L(M_{11})$; \ac{d.h.} es \ac{ex.} eine \ac{akzept.} Berechnung von $M$ f�r $w$. Da $q_2$ einziger \ac{akzept.} Zustand, muss \ac{Ber.} in $q_0$ starten und in $q_2$ enden. Der einzige Pfad von $q_0$ zu $q_2$ f�hrt �ber $q_1$:
		\[(q_0, w) \vd^* (q_1, u) \vd^* (q_2, \epsilon)\]
		Jede \ac{Berechn.} von $M$ kann h�chstens einmal in eine Konfiguration mit Zustand $q_1$ gelangen, denn $q_1 \notin \delta(q_1, a)$ f�r $a \in \Sigma_{\tx{Bool}}$. Die Berechnung muss daher wie folgt aussehen:
		\[(q_0, w) \vd^* (q_0, abu) \vd (q_1, bu), \vd (q_2, u) \vd^* (q_2, \epsilon)\]
		mit $a, b \in \SigB$. Die einzige M�glichkeit zu $q_1$ zu gelangen, ist die \ac{Anw.} der Transition $q_0 \stack{1}{\ra} q_1$, und somit folgt $a = 1$. Die einzige \ac{M�gl.} $q_1$ zu verlassen ist $q_1 \stack{1}{\ra} q_2$; daher $b = 1$ gelten, also insgesamt
		\[(q_0, w) \vd^* (q_0, 11u) \vd (q_1, 1u) \vd (q_2, u) \vd^* (q_2, \epsilon)\]
		\Ra offenbar hat $w$ $11$ enthalten, also $w \in L$
\end{itemize}
\Mark{Lemma 3.2}
\end{lemma}

\section{\texorpdfstring{$\epsilon$-Nichtdeterministische Automaten}{Epsilon-Nichtdeterministische Automaten}}
$\epsilon$-Transitionen\\
\begin{center}
\begin{tikzpicture}[->,>=stealth',shorten >=1pt,auto,node distance=2.8cm,semithick]
\node[state,initial,initial text=] (n0) {};
\node[state] (n1) [right of=n0] {};
\node[state] (n2) [right of=n1] {};
\node[state,accepting] (n3) [right of=n2] {};

\path[->] (n0) edge[loop below] node[right] {$1, 0$} ();
\path[->] (n3) edge[loop below] node[right] {$1, 0$} ();
\path[->] (n0) edge node[above] {$\epsilon$} (n1);
\path[->] (n1) edge node[above] {$1$} (n2);
\path[->] (n2) edge node[below] {$1$} (n3);
\path[->] (n2) edge[bend left] node[above] {$\epsilon$} (n3);
\end{tikzpicture}
\end{center}
\Img{GI-28.04.2009-IMG-1}

\subsection{\texorpdfstring{NEAs mit $\epsilon$-Transitionen}{NEAs mit Epsilon-Transitionen}}
\begin{fdefinition}
Ein $\epsilon$-\ac{NEA} $\mu$ ist ein Quintupel $(Q, \Sigma, \delta, q_0, F)$.
\begin{itemize}
\item $Q, \Sigma, q, F$ wie \ac{NEA}
\item $\delta$: $Q \times (\Sigma \cup \gklamm{\epsilon}) \ra \mathcal{P}(Q)$
\end{itemize} \qed
\end{fdefinition}

\begin{notation}
$q \in \delta (q, a)$ schreiben wir als $p \stack{a}{\ra} q$ $a \in \Sigma \cup \gklamm{\epsilon}$
\end{notation}

�bernehme Begriff \ggq{Konfiguration von \ac{NEA}}. Zur Definition der Schrittrelation; an �bergang
\[(p, w) \vd (q, w)\]
\begin{center}
\begin{tikzpicture}[->,>=stealth',shorten >=1pt,auto,node distance=2.8cm,semithick]
\node[state] (n0) {$p$};
\node[state] (n1) [right of=n0] {$p_1$};
\node[state] (n2) [right of=n1] {\dots};
\node[state] (n3) [right of=n2] {$p_4$};
\node[state] (n4) [right of=n3] {$q$};

\path[->] (n0) edge node[above] {$\epsilon$} (n1);
\path[->] (n1) edge node[above] {$\epsilon$} (n2);
\path[->] (n2) edge node[above] {$\epsilon$} (n3);
\path[->] (n3) edge node[above] {$\epsilon$} (n4);
\end{tikzpicture}
\end{center}
\Img{GI-28.04.2009-IMG-2}
nicht zu erkennen, wie viele $\epsilon$-Transitionen benutzt wurden. \Ra Erlaube derartige �berg�nge mit \ac{bel.} vielen $\epsilon$-Transitionen.

\begin{fdefinition}
Sei $M = (Q, \Sigma, \delta, q_0, F)$ ein $\epsilon$-\ac{NEA}. Die $\epsilon$-H�lle $\epsilon_c$: $Q \ra \mathcal{P}(Q)$ ist \ac{def.} durch
\[\epsilon_c (q) = \gklamm{q} \cup \bigcup_{q' \in \delta(q, \epsilon)} \epsilon_c(q')\]
\Mark{Definition 3.13}
\end{fdefinition}

\begin{fdefinition}
Sei $M = (Q, \Sigma, \delta, q_0, F)$ ein $\epsilon$-\ac{NEA}. Ein Konfigurations�bergang (ein Schritt) ist eine Relation $\vd:~(Q \times \Sigma^*) \times (Q \times \Sigma^*)$, die \ac{def.} ist durch
\[(q, aw) \vd (p, w) \tx{ falls } p \in \delta(q, a)\]
und
\[(q, w) \vd (p, w) \tx{ falls } p \in \epsilon_c (q)\]
wobei $a \in \Sigma, w \in \Sigma^*$ und $p, q \in Q$.
\Mark{Definition 3.14}
\end{fdefinition}

\begin{fdefinition}
f�r Berechnungen und von $\epsilon$-\ac{NEA} \ac{akzept.} Sprache analog zu \ac{NEA}.
\end{fdefinition}

\subsection{\texorpdfstring{�quivalenz von $\epsilon$-NEA und NEA}{�quivalenz von Epsilon-NEA und NEA}}
\begin{center}
\begin{tikzpicture}[->,>=stealth',shorten >=1pt,auto,node distance=2.8cm,semithick]
\node[state] (n0) {$p$};
\node[state] (n1) [right of=n0] {$q$};
\node[state] (n2) [right of=n1] {$q'$};
\node[state] (n3) [right of=n2] {$r$};

\path[->] (n0) edge node[above] {$\epsilon$} (n1);
\path[->] (n1) edge node[above] {$a$} (n2);
\path[->] (n2) edge node[above] {$\epsilon$} (n3);
\path[->] (n0) edge[bend right] node[below] {$a$} (n3);
\end{tikzpicture}
\end{center}
\Img{GI-28.04.2009-IMG-3}

\begin{beispiel}
\mbox{}\par
\begin{center}
\begin{tikzpicture}[->,>=stealth',shorten >=1pt,auto,node distance=2.8cm,semithick]
\node[state,initial,initial text=] (n0) {$q_0$};
\node[state] (n1) [right of=n0] {$q_1$};
\node[state] (n4) [below right of=n1] {$q_4$};
\node[state] (n2) [above right of=n4] {$q_2$};
\node[state,accepting] (n3) [right of=n2] {$q_3$};
\node[state] (n5) [below of=n4] {$q_5$};

\path[->] (n1) edge[loop below] node[right] {$0$} ();
\path[->] (n1) edge node[above] {$0$} (n2);
\path[->] (n4) edge node[above] {$0$} (n3);
\path[->] (n4) edge[bend right] node[left] {$1$} (n5);
\path[->] (n5) edge[bend right] node[right] {$1$} (n4);

\path[->,color=red] (n0) edge node[color=red, above] {$\epsilon$} (n1);
\path[->,color=red] (n2) edge node[color=red, above] {$\epsilon$} (n3);
\path[->,color=red] (n0) edge[bend right] node[color=red, above] {$\epsilon$} (n4);

\path[->,color=green] (n0) edge[bend left] node[color=green, below right] {$0$} (n1);
\path[->,color=green] (n1) edge[bend left] node[color=green, below] {$0$} (n3);
\path[->,color=green] (n0) edge[bend right] node[color=green, above] {$1$} (n5);
\path[->,color=green] (n0) edge[bend left] node[color=green, below] {$0$} (n2);
\path[->,color=green] (n0) edge[bend left] node[color=green, above] {$0$} (n3);
\end{tikzpicture}
\end{center}
\Img{GI-28.04.2009-IMG-4}
\end{beispiel}

\begin{fsatz}
Zu jedem $\epsilon$-\ac{NEA} l�sst sich \textbf{effektiv} ein \textbf{�quivalenter} NEA konstruieren.
\Mark{Satz 3.1}
\end{fsatz}

\begin{fbeweis}
Sei $M_{\epsilon} = (Q, \Sigma, \delta_{\epsilon}, q_0, F_{\epsilon})$ ein $\epsilon$-\ac{NEA}. Konstruktionsidee f�r \ac{NEA}: Eliminiere $\epsilon$-Transitionen und kompensiere durch Hinzunahme neuer Transitionen.\\
Formal: Konstruiere NEA $M = (Q, \Sigma, \underline{\delta}, q_0, \underline{F})$ aus $M_{\epsilon}$ wie folgt:
\[\textcolor{green}{r \in \delta(p, a)} \Lra \exists q, q' \in Q:~q \in \epsilon_c (p) \und r \in \epsilon_c (q') \und q' \in \delta (q, a)\]

\begin{center}
\begin{tikzpicture}[->,>=stealth',shorten >=1pt,auto,node distance=2.8cm,semithick]
\node[state] (n0) {$p$};
\node[state] (n1) [right of=n0] {$q$};
\node[state] (n2) [right of=n1] {$q'$};
\node[state] (n3) [right of=n2] {$r$};

\path[->, color=red] (n0) edge node[above, color=red] {$\epsilon$} (n1);
\path[->] (n1) edge node[above] {$a$} (n2);
\path[->, color=red] (n2) edge node[above, color=red] {$\epsilon$} (n3);
\path[->, color=green] (n0) edge[bend right] node[below, color=green] {$a$} (n3);
\end{tikzpicture}

\begin{tikzpicture}[->,>=stealth',shorten >=1pt,auto,node distance=2.8cm,semithick]
\node[state,initial,initial text=] (n0) {};
\node[state,accepting] (n1) [right of=n0] {};
\path[->] (n0) edge node[above] {$\epsilon^*$} (n1);
\end{tikzpicture}
\end{center}
und
\[F = \begin{cases}
F_{\epsilon} \cup \gklamm{q_0} & \tx{falls } \epsilon_c (q_0) \cap F_{\epsilon} \neq \emptyset\\
F_{\epsilon} & \tx{sonstige}
\end{cases}\]
Effektive Berechnung: ohne Beweis.

\begin{enumerate}[label=\alph*)]
\item Bleibt zu zeigen $L(M) = L(M_2)$.\\
		$L(M_{\epsilon}) \subseteq L(M)$. Sei $w = a_1 \dots a_n \in L(M_\epsilon)$ und $q_0 \stack{w}{\ra} p_n$ (also: $w$ ist Pfadbeschreibung von $q_0$ nach $p_n$ \ac{incl.} der $\epsilon$-Transitionen). \ac{bzw.}
		\[(q_0, a_1 \dots a_n) \vd_{M_{\epsilon}} (p_1, a_2 \dots a_n) \vd_{M_{\epsilon}} \dots \vd_{M_{\epsilon}} (p_{n - 1}, a_n) \vd (p_n, \epsilon)\]
		ist eine \ac{akzept.} Berechnung von $M_{\epsilon}$ f�r $w$.
		
		Wir betrachten Indizes $i_1, \dots, i_m$ mit $a_{ij} \neq \epsilon$ ($1 \klgl j \klgl m$). Offenbar ist $a_{i1} a_{i2} \dots a_{im} = w$.
		\begin{enumerate}[label=\roman*)]
		\item $(m = 0)$ es gibt keinen solchen Index $i_m$ mit $a_{ij} \neq \epsilon$, dann ist $w = \epsilon$ und $p_n \in \epsilon_c (q_0)$ sowie $p_n \in F_{\epsilon}$. $F$ wurde f�r diesen Fall aber so konstruiert, dass $q_0 \in F$.\\
				\Ra $\epsilon \in L(M)$
		
		\item $(m > 0)$ Gem�� \ac{Def.} von $\delta$ \Todo{ref zur definition} ist
				\[\tx{\ac{NEA}: } (q_0, a_{i1} \dots a_{im}) \vd (p_{i1}, a_{i2} a_{i3} \dots a_{im}) \vd \dots \vd (p_{i_{m - 1}} a_{i_m}) \vd (p_n, \epsilon)\]
				Machen wir uns das am \ac{Bsp.} des ersten Schritts nochmals klar:
				\[(q_0, a_{i_1} \dots a_{i_m}) \vd (p_{i_1}, a_{i_2} \dots a_{i_m})\]
				Angenommen es gibt im $\epsilon$-\ac{NEA} \textbf{keine} Transition
				\[q_0 \stack{q_{i_1}}{\ra} q_{i_1}\]
				\ac{d.h.} um \ac{NEA} wurde eine hinzugef�gte Kante benutzt. Offenbar ist dann $i_2 - i_1 > 1$; \ac{d.h.} im $\epsilon$-\ac{NEA} wurden $\epsilon$-Transitionen benutzt. Dann muss es laut Konstruktion $\space^*$ von $\delta$ ein $q'$ mit $q' \in \epsilon_c (q_0)$ geben, so dass $\delta '' \in \delta_{\epsilon} (q', a)$ und $q_{i_1} \in \epsilon_c (q'')$:
				\begin{center}
				\begin{tikzpicture}[->,>=stealth',shorten >=1pt,auto,node distance=2.8cm,semithick]
				\node (n0) {$q_0$};
				\node (n1) [right of=n0] {$q'$};
				\node (n2) [right of=n1] {$q''$};
				\node (n3) [right of=n2] {$p_{i_1}$};
				
				\path[->] (n0) edge node[above] {$\epsilon$} (n1);
				\path[->] (n1) edge node[above] {$a$} (n2);
				\path[->] (n2) edge node[above] {$\epsilon$} (n3);
				\path[->] (n0) edge[bend left] node[above] {$a$} (n3);
				\end{tikzpicture}
				\end{center}
				\Img{GI-29.04.2009-IMG-1}

				sonst h�tte $\epsilon$-\ac{NEA} Konfiguration ($p_{i_1}, a_{i_2} \dots a_{i_m}$) niemand erreicht. In der Konstruktion f�r $\delta$ haben wir f�r solche F�lle vereinbart, dass $p_{i_1} \in \delta(q_0, a_{i_1})$.
		\end{enumerate}

\item $L(M) \subseteq L(M_\epsilon)$: Sei $w = a_1 \dots a_n \in L(M)$ und
		\[q_0, a_1 \dots a_n) \vd_{\mu} (p_1, a_2 \dots a_n) \vd_{\mu} \dots \vd_{\mu} (p_n, \epsilon)\]
		mit $p_n \in F$ eine \ac{akzept.} Berechnung des \ac{NEA} $M$. Gem�� \ac{Def.} von $\delta$ \Todo{ref auf definition} gilt nun aber auch, dass
		\[(q_0, a_1 \dots a_n) \vd_{M_\epsilon}^* (p_1, a_2 \dots a_n) \vd_{M_\epsilon}^* \dots \vd_{M_\epsilon}^* (p_n, \epsilon)\]
		eine \ac{akzept.} Berechnung des $\epsilon$-\ac{NEA} ist. Zwei F�lle sind zu unterscheiden:
		\begin{enumerate}[label=\roman*)]
		\item $p_n \in F_{\epsilon}$ folgt sofort $w \in L(M_{\epsilon})$
		\item $p_n \in F - F_{\epsilon}$; \ac{d.h.} $p_n = q_0$: Gem�� \ac{Def.} von $F$ \Todo{ref zur definition} folgt nun f�r ein $\hat{p} \in F_{\epsilon}$: $\hat{p} \in \epsilon_c (q_0)$ und somit \ac{ex.} eine Berechnung
				\[(q_0, a_1 \dots a_n) \vd_{\mu}^* (p_1, a_2 \dots a_n) \vd_{\mu}^* \dots \vd_{\mu}^* (\hat{p}, \epsilon)\]
		\end{enumerate}
\end{enumerate} \qed
\end{fbeweis}

$L = \gklamm{w \in \SigBs \vert \tx{Zahlenwert}(w) \mod 3 = 0 \und w = u111v \tx{ f�r } u, v \in \SigBs}$\\
\ac{DEA} $M$ f�r $L$ gesucht!\\
$L = \ub{\gklamm{w \in \SigBs \vert \tx{Zahlenwert}(w) \mod 3 = 0}}{L_1} \cap \ub{\gklamm{w \in \SigBs \vert w = u111v}}{L_2}$
\begin{center}
\begin{tikzpicture}[->,>=stealth',shorten >=1pt,auto,node distance=2.8cm,semithick]
\node[state, initial, initial text=, accepting] (n0) {$\overline{0}$};
\node[state] (n1) [right of=n0] {$\overline{1}$};
\node[state] (n2) [right of=n1] {$\overline{2}$};
\node [left of=n0, node distance=1.5cm] {$M_1$:};

\path[->] (n0) edge[loop above] node[right] {$0$} ();
\path[->] (n2) edge[loop above] node[right] {$1$} ();

\path[->] (n0) edge[bend left] node[above] {$1$} (n1);
\path[->] (n1) edge[bend left] node[below] {$1$} (n0);
\path[->] (n1) edge[bend left] node[above] {$0$} (n2);
\path[->] (n2) edge[bend left] node[below] {$0$} (n1);
\end{tikzpicture}
\end{center}
\Img{GI-29.04.2009-IMG-2.1}

\begin{center}
\begin{tikzpicture}[->,>=stealth',shorten >=1pt,auto,node distance=2.8cm,semithick]
\node[state, initial, initial text=] (n0) {$\epsilon$};
\node[state] (n1) [right of=n0] {$1$};
\node[state] (n2) [right of=n1] {$11$};
\node[state, accepting] (n3) [right of=n2] {$111$};
\node [left of=n0, node distance=1.5cm] {$M_2$:};

\path[->] (n0) edge[loop above] node[right] {$0$} ();
\path[->] (n3) edge[loop above] node[right] {$0,1$} ();

\path[->] (n0) edge node[above] {$1$} (n1);
\path[->] (n1) edge node[above] {$1$} (n2);
\path[->] (n2) edge node[above] {$1$} (n3);

\path[->] (n1) edge[bend left] node[below] {$0$} (n0);
\path[->] (n2) edge[bend right] node[above] {$0$} (n0);
\end{tikzpicture}
\end{center}
\Img{Gi-29.04.2009-IMG-2.2}

\begin{center}
\begin{tikzpicture}[->,>=stealth',shorten >=1pt,auto,node distance=2.8cm,semithick]
\node[state, initial, initial text=] (n0) {$\begin{array}{c}\overline{0}\\\epsilon\end{array}$};
\node[state] (n1) [right of=n0] {$\begin{array}{c}\overline{1}\\1\end{array}$};
\node[state] (n2) [right of=n1] {$\begin{array}{c}\overline{1}\\1\end{array}$};
\node[state] (n3) [right of=n2] {$\begin{array}{c}\overline{2}\\\epsilon\end{array}$};
\node[state] (n4) [below of=n1] {$\begin{array}{c}\overline{0}\\11\end{array}$};
\node[state] (n5) [below of=n2] {$\begin{array}{c}\overline{2}\\1\end{array}$};
\node (n6) [right of=n5] {$\dots$};
\node[state, accepting] (n7) [right of=n6] {$\begin{array}{c}\overline{0}\\111\end{array}$};
\node[left of=n0,node distance=2.0cm] {$M$:};

\path[->] (n0) edge[loop above] node[right] {$0$} ();

\path[->] (n0) edge node[above] {$1$} (n1);
\path[->] (n1) edge node[above] {$0$} (n2);
\path[->] (n2) edge node[above] {$0$} (n3);
\path[->] (n1) edge node[right] {$1$} (n4);
\path[->] (n2) edge node[right] {$1$} (n5);
\end{tikzpicture}
\end{center}
\Img{GI-29.04.2009-IMG-3}

\subsection{�quivalenz von NEA und DEA}
\begin{center}
\begin{tikzpicture}[->,>=stealth',shorten >=1pt,auto,node distance=2.8cm,semithick]
\node (n0) {$\epsilon$-NEA};
\node (n1) [right of=n0] {NEA};
\node (n2) [right of=n1] {DEA};

\path[->] (n0) edge[bend left] (n1);
\path[->] (n1) edge[bend left] (n2);
\path[->] (n0) edge[bend right] (n2);
\path[->] (n1) edge (n0);
\path[->] (n2) edge (n1);
\end{tikzpicture}
\end{center}
\Img{GI-30.04.2009-IMG-1}

\begin{center}
\begin{tikzpicture}[->,>=stealth',shorten >=1pt,auto,node distance=2.8cm,semithick]
\node[state, initial, initial text=] (n0) {$A$};
\node[state] (n1) [right of=n0] {$B$};
\node[state,accepting] (n2) [right of=n1] {$C$};
\node at (n2.east) [right,node distance=0.5cm] {$Q = \{A, B, C\}$};

\path[->] (n0) edge[loop above] node[right] {$1, 0$} ();
\path[->] (n0) edge node[above] {$0$} (n1);
\path[->] (n1) edge node[above] {$0$} (n2);
\end{tikzpicture}
\end{center}
\Img{GI-30.04.2009-IMG-2.1}

\begin{center}
\begin{tikzpicture}[->,>=stealth',shorten >=1pt,auto,node distance=2.8cm,semithick]
\node[state, initial, initial text=] (n0) {$\{A\}$};
\node[state] (n1) [right of=n0] {$\{A, B\}$};
\node[state,accepting] (n2) [right of=n1] {$\{A, B, C\}$};
\node at (n2.east) [right,node distance=0.5cm] {$Q_P = \mathcal{P} (Q)$};

\path[->] (n0) edge[loop above] node[right] {$1, 0$} ();
\path[->] (n0) edge[loop below] node[right] {$0$} ();
\path[->] (n0) edge node[below] {$0$} (n1);
\path[->] (n1) edge node[above] {$0$} (n2);
\path[->] (n2) edge[bend right] node[above] {$1$} (n0);
\path[->] (n1) edge[bend right] node[below] {$1$} (n0);
\end{tikzpicture}
\end{center}
\Img{GI-30.04.2009-IMG-2.2}

$L = \gklamm{w \in \SigBs \vert w^R \tx{ enth�lt an $n$-ten Stelle eine } 1}$ $Q_D = \mathcal{P}(Q)$
\begin{center}
\begin{tikzpicture}[->,>=stealth',shorten >=1pt,auto,node distance=2.8cm,semithick]
\node[state, initial, initial text=] (n0) {$S$};
\node[state] (n1) [right of=n0] {$A$};
\node[state] (n2) [right of=n1] {$T_1$};
\node[state, accepting] (n3) [right of=n2] {$T_2$};

\path[->] (n0) edge[loop above] node[right] {$1,0$} ();
\path[->] (n0) edge node[above] {$1$} (n1);
\path[->] (n1) edge node[above] {$1,0$} (n2);
\path[->] (n2) edge node[above] {$1,0$} (n3);
\end{tikzpicture}
\end{center}
\Img{GI-30.04.2009-IMG-3.1}

\begin{center}
\begin{tikzpicture}[->,>=stealth',shorten >=1pt,auto,node distance=2.8cm,semithick]
\node[state,initial, initial text=] (n0) {$\{S\}$};
\node[state] (n1) [right of=n0] {$\{S, A\}$};
\node[state] (n2) [right of=n1] {$\{S, A, T_1\}$};
\node[state,accepting] (n3) [right of=n2] {$\{S, A, T_1, T_2\}$};
\node[state] (n4) [below of=n1] {$\{S, T_1\}$};
\node[state] (n5) [below of=n2] {$\{S, A, T_2\}$};
\node [right of=n5] {$\dots$};

\tikzloop{n0}{0}{above}{right}

\tikzedge{n0}{n1}{1}{}{above}
\tikzedge{n1}{n2}{1}{}{above}
\tikzedge{n2}{n3}{1}{}{above}
\tikzedge{n1}{n4}{0}{}{right}
\tikzedge{n4}{n5}{1}{}{above}
\end{tikzpicture}
\end{center}
\Img{GI-30.04.2009-IMG-3.2}

\begin{fdefinition}[Potenzmengenkonstruktion]
Sei $M_N = \rklamm{Q_N, \Sigma, \delta_N, q_0, F_N}$ ein \ac{NEA}. WIr konstruieren aus $M_N$ einen \ac{DEA} $M_D = \rklamm{Q_D, \Sigma, \delta_D, \rklamm{q_0}, F_D}$ wie folgt:
\begin{itemize}
\item $Q_D = \mathcal{P}(Q_N)$
\item $F_D = \gklamm{S \in Q_D \vert S \cap F_N \neq \emptyset}$
\item F�r jedes $S \in Q_D$ und jedes Eingabesymbol $a \in \Sigma$
		\[\delta_D (S, a) = \bigcup_{q \in S} \delta_N (q, a)\]
\end{itemize}
Man sagt $M_D$ entsteht aus $M_N$ durch Anwendung der \ac{PMK}.
\Mark{Definition 3.15}
\end{fdefinition}

\begin{beispiel}
Gegeben Sei \ac{NEA} $M_N = \rklamm{\rklamm{A, B, C}, \SigB, \delta_N, A, \rklamm{B}}$ mit\\
\begin{multicols}{2}
$\delta_N (A, 0) = \rklamm{C}$\\
$\delta_N (B, 0) = \rklamm{B}$\\
$\delta_N (C, 0) = \emptyset$\\
$\delta_N (A, 1) = \rklamm{A, B}$\\
$\delta_N (B, 1) = \rklamm{C}$\\
$\delta_N (C, 1) = \emptyset$
\end{multicols}
\Insert{GI-05.05.2009-IMG-1}
Einen \ac{DEA} $M_D = \rklamm{\mathcal{P} \rklamm{\gklamm{A, B, C}}, \SigB, \delta_D \gklamm{A}, F_D}$ zu $M_N$ erhalten wir durch \ac{Anw.} der \ac{PMK} wie folgt:
\begin{itemize}
\item $Q_D = \gklamm{\emptyset, \gklamm{A}, \underline{\gklamm{B}}, \gklamm{C}, \gklamm{A, \underline{B}}, \gklamm{A, C}, \gklamm{\underline{B}, C}, \gklamm{A, \underline{B}, C}}$
\item $F_D = \gklamm{\gklamm{B}, \gklamm{A, B}, \gklamm{B, C}, \gklamm{A, B, C}}$ denn $F_N = \gklamm{B}$.
\item $\delta_D \rklamm{\emptyset, 0} = \emptyset$ $\delta_D \rklamm{\emptyset, 1} = \emptyset$ (Dead State)
\item $\delta_D \rklamm{\gklamm{A}, 0} = \gklamm{C}$ $\delta_D \rklamm{\gklamm{A}, 1} = \gklamm{A, B} = \delta_D \rklamm{A, 1}$
\item $\delta_D \rklamm{\gklamm{B}, 0} = \gklamm{B}$ $\delta_D \rklamm{\gklamm{B}, 1} = \gklamm{C}$\\
		$\delta_D \rklamm{\gklamm{C}, 0} = \emptyset$ $\delta_D \rklamm{\gklamm{C}, 1} = \emptyset$\\
		$\delta_D \rklamm{\gklamm{A, B}, 0} = \bigcup_{q \in \gklamm{A, B}} \delta_N (q, 0) = \delta_N (A, 0) \cup \delta_N (B, 0) = \gklamm{C} \cup \gklamm{B} = \gklamm{B, C}$\\
		$\delta_D \rklamm{\gklamm{A, B}, 1} = \gklamm{A, B, C}$\\
		$\vdots$
\end{itemize}
\Insert{GI-05.05.2009-IMG-2}
\Mark{Beispiel 3.4}
\end{beispiel}

\begin{fsatz}
Sei $M_N = \rklamm{Q_N, \Sigma, \delta_N, q_0, F_N}$ ein \ac{NEA} und $M_D = \rklamm{Q_D, \Sigma, \delta_D, \gklamm{q_0}, F_N}$ ein \ac{DEA} der durch \ac{PMK} aus $M_N$ entsteht. Dann gilt $L(M_N) = L(M_D)$.
\Mark{Satz 3.2}
\end{fsatz}

\begin{fbeweis}
Wir beweisen zun�chst
\[\hat{\delta}_D \rklamm{\gklamm{q_0}, w} = \hat{\delta}_N \rklamm{q_0, w} \tx{\textcircled{$\star$}}\]
denn dann folgt:
\begin{align*}
\underline{w \in L\rklamm{M_N}} &\Lra \exists q_f \in F_N:~\rklamm{q_0, w} \vd_{M_N}^* \rklamm{q_f, \epsilon}\\
&\Lra \exists q_f \in F_N:~q_f \in \hat{\delta}_N \rklamm{q_0, w}\\
&\Lra \exists q_f \in F_N:~q_f \in \hat{\delta}_D \rklamm{\gklamm{q_0}, w}\\
&\Lra \hat{\delta}_D \rklamm{\gklamm{q_0}, w} \cap F_N \neq \emptyset\\
&\Lra \hat{\delta}_D \in F_D\\
&\Lra \underline{w \in L(M_D)}
\end{align*}
Wir beweisen \textcircled{$\star$} durch \ac{vollst.} Induktion �ber die Wortl�nge von $w$.
Induktionsanfang: $w = \epsilon$ \ac{bzw.} $\betrag{w} = 0$.\\
Laut \ac{Def.} von $\hat{\delta}_D$ \Todo{Ref to Definition} \ac{bzw.} $\hat{\delta}_N$ gilt:
\[\hat{\delta}_N \rklamm{q_0, \epsilon} = \gklamm{q_0} = \hat{\delta}_D \rklamm{q_0, \epsilon}\]
Induktionsschritt: Es gelte $\betrag{w} = n + 1$ \ac{bzw.} $w = \underline{u}a$. Laut Induktionsvoraussetzung gilt
\[\hat{\delta}_N \rklamm{q_0, u} = \hat{\delta}_D \rklamm{\gklamm{q_0}, u} = \gklamm{q_1, q_2, \dots, q_R}\]
Aus der induktiven Definition von $\hat{\delta}_N$ wissen wir
\[\hat{\delta}_N \rklamm{q_0, ua} = \bigcup_{q \in \hat{\delta}_N (q_0, u)} \delta_N (q, a) = \bigcup_{i = 1}^k \delta_N (q_i, a)\]
Die \ac{PMK} liefert uns andererseits
\[\hat{\delta}_D \rklamm{\gklamm{q_1, \dots, q_k}, a} = \bigcup_{i = 1}^k \delta_N \rklamm{q_i, a}\]
Jetzt nutzen wir die \ac{IV} und die \ac{indukt.} \ac{Def.} von $\hat{\delta}_D$ aus und erhalten
\begin{align*}
\hat{\delta}_D \rklamm{\gklamm{q_0}, ua} &\ubl[8mm]{=}{\ac{Def.} $\hat{\delta}_D$} \delta_D \rklamm{\hat{\delta}_D \rklamm{\gklamm{q_0}, u}, a} \ubl[8mm]{=}{\ac{IV}} \delta_D \rklamm{\hat{\delta}_D \rklamm{q_0, u}, a}\\
&= \delta_D \rklamm{\gklamm{q_1, \dots, q_R}, a} = \ubl[8mm]{=}{\ac{PMK}} \bigcup_{i = 1}^k \delta_N \rklamm{q_i, a} \ubl[8mm]{=}{\ac{Def.} $\hat{\delta}_N$} \hat{\delta}_N \rklamm{q_0, uq} = \hat{\delta}_N \rklamm{q_0, w}
\end{align*}
\end{fbeweis}

\subsection{\texorpdfstring{Ausdehnung auf $\epsilon$-NEAs}{Ausdehnung auf Epsilon-NEAs}}
\ggq{Werkzeug}: $\epsilon_c$ ($\epsilon$-H�lle)\\
Neu:
\begin{itemize}
\item $\delta_D \rklamm{S, a} = \gklamm{q \in Q_N \vert q \in \epsilon_c \rklamm{\delta_N \rklamm{r, a}} \tx{ f�r ein} r \in Q_N}$
\item $q_{0_D} = \epsilon_c \rklamm{q_{0_N}}$
\end{itemize}

\section{Minimierung von Automaten}
\Mark{Section 3.4}
\Insert{GI-06.05.2009-IMG-1}
Minimieren von Automaten bedeutet
\begin{itemize}
\item entfernen unerreichbarer Zust�nde \Todo{G und H mit rot durchstreichen in IMG-1}
\item Zusammenfassen �quivalenter Zust�nde
\end{itemize}
\begin{fdefinition}
Zwei Zust�nde $p$, $q$ hei�en �quivalent \ac{gdw.}
\[\forall w \in \Sigma^*:~\hat{\delta} (p, w) \in F \Lra \hat{\delta} (q, w) \in F\]
(in Zeichen $p \equiv q$)
\end{fdefinition}

Problem $\Sigma^*$ enth�lt unendlich viele W�rter.\\
Muss man wirklich alle pr�fen?\\
Nein denn ein Automat mit $\betrag{Q} = n$ Zust�nden durchl�uft mindestens einen Zustand zweimal, wenn er ein Wort $\betrag{w} \grgl n$ abarbeitet.\\
Wie aufw�ndig ist dieses Verfahren? Wir m�ssen alle Zustandspaare untersuchen, das sind $\binom{n}{2}$.\\
F�r jedes Paar m�ssen alle W�rter der L�nge $n$ getestet werden.

\begin{beispiel}
$\SigB$
\begin{align*}
T (n) = \binom{n}{2} \mal \sum_{i = 0}^{n - 1} 2^i &= \frac{n!}{(n - 2)! \mal 2!} \mal 2^{n} - 1\\
&= \rklamm{n^2 - n} \mal \rklamm{2^{n - 1} - \frac{1}{2}}
\end{align*}
\begin{center}
\begin{tabular}{c|c}
$n$ & $T(n)$ ca.\\
\hline
$10$ & $46.000$\\
$11$ & $112.000$\\
\vdots & \vdots\\
$30$ & $467.077.00.000$
\end{tabular}
\end{center}
\end{beispiel}

\subsection{Eigenschaften einer �quivalenzrelation}
\begin{itemize}
\item $q \equiv q$ (Reflexivit�t)
\item $q \equiv p \und p \equiv r \Ra q \equiv r$ (Transitivit�t)
\item $q \equiv p \Lra \equiv q$ (Symmetrisch)
\end{itemize}
Angenommen $p \equiv q$, dann gilt
\[\forall w \in \Sigma^*:~\hat{\delta} (p, w) \in F \Lra \hat{\delta} (q, w) \in F\]
Betrachten wir zwei Berechnungen f�r konkretes $w = a_1 \dots a_n$:
\begin{align*}
&(p, a_1 \dots a_n) \vd (p_1, a_2 \dots a_n) \vd \dots \vd (p_n, \epsilon)\\
&(q, a_1 \dots a_n) \vd (q_1, a_2 \dots a_n) \vd \dots \vd (q_n, \epsilon)
\end{align*}
Offenbar gilt $p_i \equiv q_i$ ($1 \klgl i \klgl n$)

\begin{beobachtung}
\mbox{}\par
\begin{enumerate}
\item $p \equiv q \Ra \forall a \in \Sigma:~\delta (p, a) \equiv \delta (q, a)$
\item $p \equiv q \Ra p \in F \Lra q \in F$
\end{enumerate}
\Insert{GI(U)-07.05.2009-IMG-1}
\Insert{GI(U)-07.05.2009-IMG-2}
\begin{center}
\begin{tabular}{c|c|c|c|c|c}
  & F						& E & D 						& C & B\\\hline
A & \tc{x}						& \tc{x} & \tc{x} 						& \tc{x} & \tc{x}\\\hline
B & \tc{x}						& \tc{x} & \tc{x} \tiny{A,D}&   &  \\\hline
C & \tc{x}						& \tc{x} & \tc{x} \tiny{A,D}&   &  \\\hline
D & \tc{x}						& \tc{x} &   						&   &  \\\hline
E & \tiny{B,C}	&   &   		&   &  \\\hline
\end{tabular}
\end{center}
\Insert{GI(U)-07.05.2009-TAB-3}
\Insert{GI(U)-07.05.2009-IMG-4}
\end{beobachtung}

%\chapter{Grammatiken}
\Mark{Chapter 4}
\section{Chomsky-Hierarchie}
\Mark{Section 4.1}
\begin{fdefinition}[Grammatik]
Eine \indexb{Grammatik $G$} ist ein Quadrupel $G = \rkl{\Sigma_N, \Sigma_T, P, S}$ aus
\begin{itemize}
\item einem Alphabet $\Sigma_N$ (\indexb{Nichtterminalsymbole})
\item einem Alphabet $\Sigma_T$ (\indexb{Terminalsymbole})
\item einem Startsymbol $S \in \Sigma_N$
\item einer endlichen Menge von \indexb{Ersetzungsregeln} (Produktionen)
		\[P \subseteq \Sigma^* \Sigma_N \Sigma^* \times \Sigma^* \tx{ mit } \Sigma^* = \Sigma_T \cup \Sigma_N\]
		statt $(\alpha, \beta) \in P$ schreiben wir $\alpha \ural{G}{} \beta$
\end{itemize}
\Mark{Definition 4.1}
\end{fdefinition}

\begin{beispiel}
\fnewline
\begin{align*}
\ubl[5em]{S}{nicht terminal Symbole} \ra \epsilon,~S \ra \ubl[5em]{O}{terminal Symbol}S,~S \ra 1S~&\Vert Ableit-Regeln\\
S \vd 1S \vd 11S \vd 11OS \vd 110~&\Vert Ableitung
\end{align*}
\end{beispiel}

\begin{notation}
\fnewline
\begin{itemize}
\item Terminalsymbole als Kleinbuchstaben $a$, $b$, $c$, $d$
\item W�rter aus Terminalsymbole $u$ $v$ $w$
\item Nichtterminalsymbole als Gro�buchstaben $A$, $B$, $C$
\item W�rter aus $(\Sigma_T \cup \Sigma_N)^*$ als griechische Buchstaben $\alpha$, $\beta$, $\gamma$
\end{itemize}

$G_2 = \rkl{\gklamm{S}, I_{\tx{Bool}}, P_2, s}$, $p_2 = \gklamm{S \ra \epsilon, s \ra 0S1}$ Vereinfacht gilt oft nur die Produktion an $S \ra \epsilon \vert 0S1$
\end{notation}

\begin{fdefinition}
Sei $G = \rkl{I_N, \Sigma_T, P, S}$ eine Grammatik und einen $\alpha, \beta \in \rkl{\Sigma_N \cup \Sigma_T}^*$
\begin{enumerate}
\item $\alpha \vd_G \beta \Lra \exists w_1, w_2 \wedge \exists \rkl{y, \delta} \in P$\\
		$a = w_1 y w_2 \wedge \beta = w_1 \delta w_2$\\
		wir sagen, dass $\beta$ aus $\alpha$ in einem Ableitungsschritt in $G$ ableitbar ist.

\item $\alpha \vd_G^h \beta \Lra \exists w_0 \dots w_n \in \Sigs$, $w_0 = \alpha$ $w_n = \beta$, so dass
		\[\forall i < n:~w_i \vd w_{i + 1}\]
		wir sagen, dass $\beta$, $w_s$, $\alpha$ in $G$ in $n$ Schritten ableitbar ist. Die Folge $w_0 \vd_G w_1 \vd_G \dots \vd_G w_n$ hei�t eine Ableitung in $G$.

\item $a \vd_G^* \beta \Lra \exists n \in N_0:~ \alpha \vd_G^n \oder \alpha = \beta$
\item Falls $S \vd_G^* w$ ($w \in \Sigma_T^*$) sagen wir, dass $w$ von $G$ erzeugt wird und $L(G) = \gklamm{w \in \Sigma_T^* \vert s \vd_G^* w}$ ist die von $G$ erzeugte Sprache.
		\[S \ra \epsilon \vert 0 \vert 1 \vert SS \rkl{L(G) = \SigBs}\]
\end{enumerate}
\Mark{Definition 4.2}
\end{fdefinition}

\begin{fdefinition}[Chomsky-Hierarchic/Sprachfassen]
\index{Chomsky-Hierarchic}\index{Sprachfassen}
Sei $G = \rkl{\Sigma_N, \Sigma_T, P, S}$ eine Grammatik
\begin{enumerate}
\item $G$ hei�t Typ-0 Grammatik \index{Typ-0 Grammatik}
\item $G$ hei�t Typ-1 Grammatik, falls jede Produktion die Form $w_1 A w_2 \ra w_1 \beta w_2$ mit $\betrag{\beta} \grgl 1$, oder $S \ra \epsilon$ hat, $A \in \Sigma_N, w_1, w_2, \beta \in \rkl{\Sigma_T \cup \Sigma_N}^*$ \Ra kontextsensitive Grammatik \index{Typ-1 Grammatik}
\item $G$ hei�t Typ-2 Grammatik (kontextfreie Grammatik), falls $\forall \rkl{\alpha, \beta} \in P:~\alpha \in \Sigma_N \und \beta \in \rkl{\Sigma_T \cup \Sigma_N}^*$ \index{Typ-2 Grammatik}
\item $G$ hei�t Typ-3 Grammatik (regul�re oder rechtslineare Grammatik) wenn $\forall \rkl{\alpha, \beta} \in P:~\alpha \in \Sigma_N \und \beta \in \rkl{\Sigma_T^* \mal \Sigma_N} \cup \Sigma_T^*$ \index{Typ-3 Grammatik}
\end{enumerate}
\Mark{Definition 4.3}
\end{fdefinition}

\section{Rechtslineare Grammatiken und regul�re Sprachen}
\Mark{Section 4.2}
Regul�re Sprachen (von \ac{DEA} erkennbar) \stack{?}{=} Typ-3 Sprachen.

\begin{lemma}
$\mathcal{L}_3$ enth�lt alle endlichen Sprachen
\Mark{Lemma 4.1}
\end{lemma}

\begin{fbeweis}
Sei $L = \gklamm{w_1, \dots, w_n}$ eine endliche Sprache �ber $\Sigma$, dann konstruieren wir eine Grammatik $G$ wie folgt
\[G = \rkl{\gklamm{S}, \Sigma, \gklamm{S \ra w_1, \dots, S \ra w_n}, S}\]
Offenbar gilt $L(G) = L$.
\end{fbeweis}

\begin{lemma}
$\mathcal{L}_3$ ist abgeschlossen bez�glich der Vereinigung.
\Mark{Lemma 4.2}
\end{lemma}

\begin{fbeweis}
Seien $L_1$ und $L_2$ zwei \ac{bel.} Sprachen aus $\mathcal{L}_3$. Dann gibt es zwei Grammatiken:
\[G_1 = \rkl{\Sigma_{N_1}, \Sigma_{T_1}, P_1, S_1}\]
und
\[G_2 = \rkl{\Sigma_{N_2}, \Sigma_{T_2}, P_2, S_2}\]
mit $L(G_1) = L_1$ und $L(G_2) = L_2$.\\
Nichtterminale k�nnen nach Bedarf umbenannt werden, daher nehmen wir \ac{o.B.d.A} an, dass $\Sigma_{N_1} \cap \Sigma_{N_2} = \emptyset$. Wir konstruieren eine Grammatik
\[G_{\cup} = \rkl{\Sigma_{N_{\cup}}, \Sigma_{T_{\cup}}, P_{\cup}, S_{\cup}}\]
mit $L(G_{\cup}) = L_1 \cup L_2$ mit
\begin{itemize}
\item $\Sigma_{N_{\cup}} = \gklamm{S_{\cup}} \cup \Sigma_{N_1} \cup \Sigma_{N_2}$
\item $\Sigma_{T_{\cup}} = \Sigma_{T_1} \cup \Sigma_{T_2}$
\item $P_{\cup} = P_1 \cup P_2 \cup \gklamm{S_{\cup} \ra S_1, S_{\cup} \ra S_2}$
\end{itemize}
Bleibt zu zeigen, dass $L(G_{\cup}) = L_{G_1} \cup L_{G_2}$
\begin{enumerate}
\item $L(G_1) \cup L(G_2) \subseteq L(G_{\cup})$. Sei $w$ ein \ac{bel.} Wort aus $L(G_1) \cup L(G_2)$. Wir unterscheiden zwei F�lle.
		\begin{enumerate}[label=\roman*)]
		\item $w \in L(G_1)$. Dann \ac{ex.} eine Ableitung
				\[S_1 \vd^* w\]
				von $w$ in $G_1$. Dann ist
				\[S_{\cup} \vd S_1 \vd^* w\]
				eine Ableitung von $w$ in $G_{\cup}$ da $G_{\cup}$ alle Produktionen aus $G_1$ enth�lt.

		\item $w \in L(G_2)$ analog.
		\end{enumerate}

\item $L(G_{\cup}) \subseteq L(G_1) \cup L(G_2)$. Seien $w \in L(G_{\cup})$ und
		\[S_{\cup} \vd \alpha_1 \vd \alpha_2 \vd \dots \vd \alpha_n = w\]
		eine Ableitung f�r $w$ in $G_{\cup}$. Wir haben nur zwei Regeln um $S_{\cup}$ zur ersetzen: $S_{\cup} = S_1$ oder $L_1 = S_2$.
		\begin{enumerate}[label=\roman*)]
		\item $\alpha_1 = S_1$, so ist
				\[S_1 \vd \alpha_2 \vd \dots \vd \alpha_n = w\]
				eine Ableitung f�r $w$ in $G_1$, da $\Sigma_{N_1} \cap \Sigma_{N_2} = \emptyset$

		\item $\alpha_1 = S_2$ analog
		\end{enumerate}
\end{enumerate}
\end{fbeweis}

\begin{lemma}
$\mathcal{L}_3$ ist abgeschlossen bez�glich Konkatenation.
\[x \ra u \Ra x \ra u S_2\]
\Mark{Lemma 4.3}
\end{lemma}

\begin{fdefinition}[Normierte regul�re Grammatik]
\index{Normierte regul�re Grammatik}
\index{Regul�re Grammatik}
\label{definition:Grammatiken_D4.4}
Eine regul�re Grammatik $G = \rkl{\Sigma_N, \Sigma_T, P, S}$ hei�t \ggq{normiert}, wenn alle Regeln der Grammatik nur eine der drei folgenden Formen hat:
\begin{itemize}
\item $S \ra \epsilon$
\item $A \ra a$
\item $B \ra bC$
\end{itemize}
f�r $A, B, C, S, \in \Sigma_{N}$; $a, b \in \Sigma_T$
\Mark{Definition 4.4}
\end{fdefinition}

\begin{lemma}
Jede regul�re Grammatik l�sst sich in eine �quivalente Grammatik $G'$ �berf�hren, die keine Kettenregeln $x \ra x$ enth�lt (wobei $x, y \in \Sigma_N$).
\end{lemma}

\begin{fbeweis}
Sei $G = \rkl{\Sigma_N, \Sigma_T, P, S}$. F�r jedes Nichtterminalsymbol $A \in \Sigma_N$ bestimmen wir die Menge $K(A)$ der NIchtterminale, die von $A$ ausgehend durch ausschlie�liche Anwendung von Kettenregeln erreicht werden k�nnen; formal
\[K(A) = \gklamm{D \in \Sigma_N \vert A \vd^* D, \tx{wobei nur Kettenregeln angewendet werden}}\]
Die Menge $K(A)$ kann algorithmisch folgenderma�en berechnet werden
\begin{align*}
K_0 (A) &= \gklamm{A}\\
K_{i + 1} (A) &= K_i (A) \cup \gklamm{B \in \Sigma_N \vert \exists C \in K_i (A) \tx{ mit } C \ra B \in P}
\end{align*}
f�r $0 \klgl i \klgl \betrag{\Sigma_N} - 2$. Dann gilt offenbar
\[K(A) = \bigcup_{i = 0}^{\un} K_i(A)\]
Die Idee ist nun, alle Ableitungsketten der Form
\[A \vd A_1 \vd A_2 \vd \dots \vd A_n \vd \alpha\]
durch Hinzunahme der Regeln $A \ra \alpha$ zu verk�rzen, wobei $\alpha \in \Sigma_r^* \cup \Sigma_r^+ \Sigma_N$\\
Formal Konstruktion: $G' = \rkl{\Sigma_N, \Sigma_T, P', S}$ mit
\[P' = \rkl{P - \Sigma_N \times \Sigma_N} \cup \bigcup_{A \in \Sigma_N} \gklamm{A \ra \alpha \vert \alpha \in \Sigma_r^* \cup \Sigma_r^+ \mal \Sigma_N \und \exists C \in K(A) \tx{ mit } C \ra \alpha}\]
\end{fbeweis}

\begin{lemma}
\label{lemma:Grammatiken_L4.5}
F�r jede regul�re Grammatik $G = \rklamm{\Sigma_N, \Sigma_T, P, S}$ existiert eine �quivalente Grammatik $G'$, die keine Regel $A \ra \epsilon$ mit $A \in \Sigma_N - \gklamm{S}$ enth�lt.
\Mark{Lemma 4.5}
\end{lemma}

\begin{fbeweis}
Wir nehmen an, dass $G$ bereits frei von Kettenregeln ist. Sei
\[Q = \gklamm{A_1 \ra \epsilon, \dots, A_R \ra \epsilon}\]
die Menge aller Regeln $P$ aus $\Sigma_N - \gklamm{S} \times \gklamm{\epsilon}$.\\
Wir konstruieren $G' = \rkl{\Sigma_N, \Sigma_r, P', S}$ mit
\[P' = \rkl{P - Q} \cup \bigcup_{i = 1}^k \gklamm{B \ra w \vert B \ra w A_i \in P}\]
Die Idee der Konstruktion besteht darin, dass die Regeln $B \ra w A_i$ und $A_i \ra \epsilon$ eine Ableitung
\[B \vd w A_i \vd w\]
in $G$ erm�glichen. Mit Einf�hrung der Regeln $B \ra w$ in $G'$ kann diese Ableitung in einem Schritt erfolgen:
\[B \vd_{G'} w\]
Wir zeigen $L(G) = L(G')$ wie �blich:
\begin{enumerate}
\item $L(G) \subseteq L(G')$. Sei $w \in L(G)$ und sei ferner
		\[S \vd_G v_1 B_1 \vd_G v_1 v_2 B_2 \vd_G \dots \vd_G v_1 \dots v_{n - 2} B_{n - 2} \vd_G v_1 \dots v_{n - 1} B_{n - 1} \vd v_1 \dots v_n = w\]
		mit $v_i \in \Sigma_T^*$ und $B_i \in \Sigma_N$ ($1 \klgl i \klgl n$) eine Ableitung von $w$ in $G$. Falls $v_n \neq \epsilon$ wurde keine Regel der Form $A \ra \epsilon$ mit $A \in \Sigma_N - \gklamm{S}$ verwendet - somit \ac{ex.} diese Ableitung auch in $G'$.\\
		Falls jedoch $v_n = \epsilon$ war die letzte Regel $B_{n - 1} \ra \epsilon_i$; dann ist die Regel (laut Konstruktion) $B_{n - 2} \ra v_{n - 1} \in P'$ und
		\[S \vd_{G'} v_1 B_1 \vd_{G'} \dots \vd_{G'} v_1 \dots v_{n - 2} B_{n - 2} \vd_{G'} v_1 \dots v_{n - 1} = w\]
		eine Ableitung f�r $w$ in $G'$.

\item $L(G') \subseteq L(G)$. Sei $w \in L(G')$ und sei
		\[S \vd_{G'}^* v_1 D \vd_{G'} v_1 v_2 = w\]
		eine Ableitung f�r $w \in G'$ mit $v_1, v_2 \in \Sigma_T^*$ und $D \in \Sigma_N$. Offenbar ist
		\[S \vd_{G}^* v_1 D\]
		eine Ableitung in $G$. F�r den letzten Ableitungsschritt gibt es zwei M�glichkeiten
		\begin{enumerate}
		\item $D \ra v_2 \in P$. Dann gilt $v_1 D \vd v_1 v_2 = w$ eine Ableitung in $G$. Somit haben wir eine Ableitung von $w$ in $G$ gefunden.
		\item $D \ra v_2 \notin P$ (Regel wurde in $G'$ hinzugef�gt), dann \ac{ex.} $A_i \in \Sigma_N$ so dass
				\[D \ra v_2 A_i, A_i \ra \Sigma \in P\]
				Dann ist
				\[S \vd_G^* v_1 D \stack{(D \ra v_2 A_i)}{\vd_G} v_1 v_2 A_i \stack{(A_i \ra \epsilon)}{\vd_G} v_1 v_2\]
				eine Ableitung f�r $w$ in $G$.
		\end{enumerate}
\end{enumerate}
Im letzten Schritt eliminieren wir Regeln der Formen $A \ra w$ und $B \ra wC$ mit $\betrag{w} \grgl 2$.
\end{fbeweis}

\begin{fsatz}
Zu jeder regul�ren Grammatik \ac{ex.} eine �quivalente normierte Grammatik.
\Mark{Satz 4.2}
\end{fsatz}

\begin{fbeweis}
Sei $G$ eine regul�re Grammatik. Wenn man die Regeln aus Definition
\vref{definition:Grammatiken_D4.4} und Lemma
\vref{lemma:Grammatiken_L4.5} \Solved{Ref zu Lemma} anwendet, erh�lt man eine
Grammatik $G_0 = \rkl{\Sigma_N, \Sigma_T, P_0, S}$ mit folgenden Eigenschaften.
\begin{enumerate}
\item $L(G_0) = L(G)$
\item $P_0$ enth�lt nur Regeln der Form
		\begin{itemize}
		\item $S \ra \epsilon$
		\item $A \ra w$ $A \in \Sigma_N$, $w \in \Sigma_T^+$
		\item $B \ra uC$ $B \in \Sigma_N$, $w \in \Sigma_T^+$
		\end{itemize}
\end{enumerate}

\begin{minipage}[t]{5cm}
\begin{align*}
A \rightarrow &a_1~\fbox{$a_2 a_3 a_4 C$}^{T_1}\\
&\downarrow\\
A \rightarrow &a_1 T_1\\
T_1 \rightarrow &a_2~\fbox{$a_3 a_4 C$}^{T_2}\\
&\downarrow\\
A \rightarrow &a_1 T_1\\
T_1 \rightarrow &a_2 T_2\\
T_2 \rightarrow &a_3~\fbox{$a_4 C$}\\
&\downarrow\\
A \rightarrow &a_1 T_1\\
T_1 \rightarrow &a_2 T_2\\
T_2 \rightarrow &a_3 T_3\\
T_3 \rightarrow &a_4 C
\end{align*}
\end{minipage}
\begin{minipage}[t]{5cm}
\begin{align*}
B \rightarrow &aC\\
B \rightarrow &b_1~\fbox{$b_2 b_3$}\\
&\downarrow^*\\
B \rightarrow & b_1 Q_1\\
Q_1 \rightarrow & b_2 Q_2\\
Q_2 \rightarrow & b_3
\end{align*}
\end{minipage}
\Img{GI-16.06.2009-IMG-1}

Um eine Grammatik in normierter Form zu erhalten m�ssen wir Regeln der Form $A \ra w$ und $B \ra uC$ mit $\betrag{w} \grgl 2$ \ac{bzw.} $\betrag{u} \grgl 2$ durch Regeln der Form $D \ra a$ und $D \ra aX$ ersetzen Idee dazu \ac{s.o.}
\end{fbeweis}

\begin{beispiel}
$S \ra aB \vert bC$\\
$B \ra b$\\
$C \ra aS \vert b$\\
$S \vd a_1 N_1 \vd a_1 a_2 N_2 \dots$

\begin{center}
\begin{tikzpicture}[->,>=stealth',shorten >=1pt,auto,node distance=2.8cm,semithick]
\node [state,initial,initial text=](S) {S};
\node [state] (B) [above right of=S] {B};
\node [state] (C) [below right of=S] {C};
\node [state,accepting](Fin) [above right of=C] {Fin};

\tikzedge{S}{B}{a}{}{};
\tikzedge{S}{C}{b}{}{};
\tikzedge{B}{Fin}{b}{}{};
\tikzedge{C}{Fin}{b}{}{};
\tikzedge{C}{S}{a}{bend left}{};
\end{tikzpicture}
\Img{GI-16.06.2009-IMG-2}
\end{center}
\end{beispiel}

\begin{fsatz}
Die Typ-3-Sprachen sind genau die Menge der regul�ren Sprachen
\[\mathcal{L}_3 = \gklamm{L \vert L \tx{ ist regul�r}}\]
\Mark{Satz 4.3}
\end{fsatz}

\begin{fbeweis}
Der Beweis erfolgt in zwei Schritten:
\begin{enumerate}
\item Es sei $L \in \mathcal{L}_3$, \ac{d.h.} f�r eine Typ-3-Grammatik $G = \rkl{\Sigma_N, \Sigma_T, P, S}$ gilt $L(G) = L$. Es gilt $w \in L(G)$ \ac{gdw.} es eine Ableitung
		\[S \vd w_1 B_1 \vd w_1 w_2 B_2 \vd \dots \vd w_1 \dots w_{n - 1} B_{n - 1} \vd w_1 \dots w_n = w\]
		f�r ein Wort $w$ in $G$ f�r Produktionen $B_{i - 1} \ra w_i B_i \in P$ mit $w_i \in \Sigma_r^*$ und $B_i \in \Sigma_N$.\\
		Wir konstruieren einen $\epsilon$-\ac{NEA} $M = \rkl{\Sigma_N \cup \Sigma_I \cup \gklamm{B_f}, \Sigma_T, \delta, S, \gklamm{B_f}}$ mit $\Sigma_N \cap \Sigma_I = \emptyset$, $\gklamm{B_f} \notsubseteq \Sigma_N \cup \Sigma_I$, der f�r jede dieser Ableitungen sukzessive die $w_i$ konsumiert. Betrachten wir Produktionen der Form
		\[B_i \ra w_j B_j\]
		Es gibt zwei M�glichkeiten:
		\begin{enumerate}
		\item $w_j = a_1 \dots a_m$ f�gen wir nicht-akzeptierende Zust�nde $T_1, \dots T_{m - 1} \in \Sigma_I$ ein und setzen
				\begin{align*}
				\delta \rkl{B_i, a_1} &= T_1\\
				\delta \rkl{T_1, a_a} &= T_2\\
				&\vdots\\
				\delta \rkl{T_{m - 1}, a_m} &= B_j
				\end{align*}

		\item $w_j = \epsilon$: wir setzen
				\begin{align*}
				\delta (B_i, \epsilon) = B_j
				\end{align*}
		\end{enumerate}
		F�r alle Produktionen der Form $B_i \ra w_j = a_1 \dots a_m$ f�gen wir neue nicht-akzeptierende Zust�nde $T_1, \dots, T_{m - 1} \in \Sigma_I$ zu $M$ hinzu und setzen
		\begin{align*}
		\delta \rkl{B_i, a_1} &= T_1\\
		\delta \rkl{T_1, a_2} &= T_2\\
		&\vdots\\
		\delta \rkl{T_{m - 1}, a_n} = B_f
		\end{align*}
		Falls $B_i \ra \epsilon$, so setzen wir
		\[\delta \rkl{B_i, \epsilon} = B_f\]
		Betrachten wir $w \in L(G)$ und die Ableitung von $w$ in $G$. Entsprechend der Konstruktion von $M$ gibt es einen Pfad von $S$ in den einzigen \ac{akzept.} ZUstand $B_f$ und somit ist $w \in L(M)$. Andersherum falls $w \in L(M)$, so wurde $M$ gerade so konstruiert, dass eine Berechnung in $M$ die Form
		\[\rkl{S, w} \vd \rkl{T_1, a_2 \dots a_m w_2 \dots w_n} \vd \rkl{T_2, a_3 \dots a_m w_2 \dots w_m} \vd \dots \vd \rkl{B_f, \epsilon}\]
		Damit ist in $G$
		\[S \vd w_1 B_1 \vd \dots \vd w_1 \dots w_n\]
		eine Ableitung von $w$ in $G$.

\item Sei $L$ eine regul�re Sprache und $M = \rkl{Q, \Sigma, \delta, q_0, F}$ ein \ac{DEA} mit $L(M) = L$ und $Q = \gklamm{q_0, \dots, q_n}$ und $\Sigma = \gklamm{a_1, \dots, a_m}$. Wir konstruieren aus $M$ eine rechtslineare Grammatik $G = \rkl{\Sigma_N, \Sigma_T, P, S}$ mit
		\begin{itemize}
		\item $\Sigma_N = Q$
		\item $\Sigma_T = \Sigma$
		\item $S = q_0$
		\end{itemize}
		f�r jede Transition $\delta \rkl{q_i, a_j} = q_k$ f�gen wir die Produktion
		\[q_i \ra a_j q_k\]
		Falls $q_k \in F$, nehmen wir zus�tzlich die Produktion
		\[q_k \ra \epsilon\]
		in $P$ auf.\\
		Korrektheit der Konstruktion trivial, selbst.
\end{enumerate}
\end{fbeweis}

%\chapter{Berechenbarkeitstheorie}
\begin{description}
\item[Algorithmus] $A: \Sigma_{\tx{Input}}^* \ra \Sigma_{\tx{Output}}^*$\\
		Algorithmus ist eine Verfahren, dass
		\begin{itemize}
		\item Eingabew�rter schrittweise verarbeitet (deterministisch).
		\item durch einen endlichen Text bis ins letzte Detail eindeutig festgelegt ist.
		\item bei Terminierung eine Ausgabe liefert, oder nicht determiniert.
		\end{itemize}
\end{description}
Wozu formales Modell, wenn wir Algorithmen auch so (\ac{z.B.} Programmiersprache) beschreiben k�nnen?
\begin{itemize}
\item geht es um Berechenbarkeit reicht intuitives Vorgehen f�r positive Antwort aus, da Programmier Sprachen eindeutige Semantik hat
\item negative Antwort ohne formales Modell nur schwer zu zeigen
\end{itemize}
Anforderungen an formales Modell
\begin{itemize}
\item einfach, um formale Beweise zu erleichtern
\item berechnungsuniversell; \ac{d.h.} alle intuitiv berechenbaren Funktionen werden erfasst.
\end{itemize}

\begin{fdefinition}
\Insert{GI-23.06.2009-IMG-1}
Eine \ac{TM} ist ein 7-Tupel
\[M = \rkl{Q, \Sigma, \Gamma, \delta, q_0, q_{\tx{accept}}, q_{\tx{reject}}}\]
mit
\begin{itemize}
\item einer endlichen Zustandsmenge $Q$
\item einem Eingabealphabet $\Sigma$, wobei das \indexb{Randsymbol} \ggq{$\cent$}\index{$\cent$ Randsymbol} und das Leerzeichen \ggq{\textvisiblespace} nicht in $\Sigma$ enthalten sind
\item einem Ausgabealphabet $\Gamma$, wobei $\Sigma \subseteq \Gamma$ und $\gklamm{\cent, w} \subseteq \Gamma$.
\item einer �bergangsfunktion
		\[\delta:~ Q - \gklamm{q_{\tx{accept}}, q_{\tx{reject}}} \times \Gamma \ra Q \times \Gamma \times \gklamm{L, R, N}\]
		mit der Eigenschaft
		\[\forall q \in Q:~\delta\rkl{q, \cent} \in Q \times \gklamm{\cent} \times \gklamm{R, N}\]
		was bedeutet das man kein $\cent$ auf das Band schreiben kann.

\item einem Startzustand $q_0 \in Q$
\item einem akzeptierenden Zustand $q_{\tx{accept}} \in Q$
\item einem verwerfenden Zustand $q_{\tx{reject}} \in Q$
\end{itemize}
$\delta(q, a) = \rkl{q_1, b, L}$ bedeutet, dass die \ac{TM} im Zustand $q$ durch Lesen des Symbols $a$ in den Zustand $q_1$ �bergehen kann, das Symbole $a$ durch $b$ ersetzt und den \ac{S-/L-Kopf} eine Position nach links bewegt.
\Mark{Definition 5.1}
\end{fdefinition}
\begin{fdefinition}
Sei $M = \rkl{Q, \Sigma, \Gamma, \delta, q_0, q_{\tx{accept}}, q_{\tx{reject}}}$ eine \ac{TM}. Eine Konfiguration $C$ von $M$ ist ein Element aus
\[\Konf{(M)} = \gklamm{\cent} \mal \Gamma^* \mal Q \mal \Gamma^+ \cup Q \mal \gklamm{\cent} \mal \Gamma^*\]
F�r ein Wort $w \in \Sigs$ ist $q_0 \cent w$ eine Startkonfiguration; eine Konfiguration der Gestalt $v_1 q_{\tx{accept}} v_2$ hei�t akzeptierende Konfiguration und eine Konfiguration der Gestalt $v_1 q_{\tx{reject}} v_2$ hei�t verwerfende Konfiguration $\rkl{v_1, v_2 \in \Gamma^*}$. Eine Konfiguration $v_1 q a v_2$ mit $v_1, v_2 \in \Gamma^*$, $a \in \Gamma$, $q \in Q$ beschreibt folgende Situation.
\begin{itemize}
\item \ac{TM} befindet sich im Zustand $q$
\item auf dem Band steht $v_1 a v_2$, wobei $v_1 = \cent v_1'$ $v_1' \in \Gamma^*$ gefolgt von unendlich vielen Blank-Symbolen \textvisiblespace
\item der \ac{S-/L-Kopf} befindet sich auf dem $\rkl{\betrag{v_1} + 1}$-ten Feld des Bandes und liefert $a$ als n�chstes Symbol.
\end{itemize}
\Mark{Definition 5.2}
\end{fdefinition}
\begin{fdefinition}
Sei $M = \rkl{Q, \Sigma, \Gamma, \delta, q_0, \qacc, \qrej}$ eine \ac{TM}. Ein Konfigurations�bergang (ein Schritt) von $M$ ist eine Relation $\vd:~\Konf{(M)} \times \Konf{(M)}$, die \ac{def.} ist durch
\begin{enumerate}
\item \ac{S-/L-Kopf} bleibt $\rkl{\delta \rkl{q, a_i} = \rkl{p, b, N}}$
		\[a_1 a_2 \dots a_{i - 1} \textcolor{red}{q} \textcolor{blue}{a_i} a_{i + 1} \dots a_n \vd_{M} a_n \dots a_{i - 1} \txcol{red}{p} \txcol{blue}{b} a_{i + 1} \dots a_n\]

\item \ac{S-/L-Kopf} geht nach links $\rkl{\delta \rkl{q, a_i} = \rkl{p, b, L}}$
		\[a_1 a_2 \dots a_{i - 1} \txcol{red}{q} \txcol{blue}{a_i} a_{i + 1} \dots a_n \vd_{M} a_1 \dots a_{i - 2} \txcol{red}{p} a_{i - 1} \txcol{blue}{b} a_{i + 1} \dots a_n\]

\item \ac{S-/L-Kopf} geht nach rechts $\rkl{\delta \rkl{q, a_i} = \rkl{p, b, R}}$
		\begin{enumerate}
		\item Kopf auf letztem Zeichen
				\[a_1 a_2 \dots a_{n - 1} \txcol{red}{q} \txcol{blue}{a_n} \vd_M a_1 a_2 \dots a_{n - 1} \txcol{blue}{b} \txcol{red}{p} \tx{\textvisiblespace}\]

		\item Kopf nicht auf letztem Zeichen
				\[a_1 a_2 \dots a_{i - 1} \txcol{red}{q} \txcol{blue}{a_i} \dots a_n \vd_M a_1 a_2 \dots a_{i - 1} \txcol{blue}{b} \txcol{red}{p} a_{i + 1} \dots a_n\]
		\end{enumerate}
		nennen wir Schrittrelation.
\end{enumerate}
\Mark{Definition 5.3}
\end{fdefinition}
\begin{fdefinition}
Sei $M = \rkl{Q, \Sigma, \Gamma, q_0, \qacc, \qrej}$ eine \ac{TM} und $w \in \Sigs$. Eine Berechnung ist eine, m�glicherweise unendliche, Kette von Konfigurationen:
\[C = C_0 C_1 \dots\]
mit $C_i \in \Konf{(M)}$ und $C_i \vd_M C_{i + 1}$ ($i \in \N$). Ist die Berechnung endlich so endet sie in einer Konfiguration $C_n = v_1 q v_2$ mit $q \in \gklamm{\qacc, \qrej}$ und $\ubl[\widthof{\tx{beginnen}}]{v_1}{muss mit $\cent$ beginnen}, v_2 \in \Gamma^*$. Eine Berechnung von $M$ hei�t
\begin{itemize}
\item akzeptierende Berechnung, wenn sie in \qacc endet.
\item verwerfende Berechnung, wenn sie in \qrej endet.
\item nicht-akzeptierende Berechnung, wenn sie in \qrej endet oder nicht terminiert.
\end{itemize}
Mit $\vd_M^*$ bezeichnen wir den reflexiven und transitiven Abschluss der Schrittrelation.
\Mark{Definition 5.4}
\end{fdefinition}
\begin{fdefinition}
Sei $M$ eine \ac{TM}. Die von $M$ akzeptierte Sprache ist
\[L(M) = \gklamm{w \in \Sigs \vert q_0 \cent w \vd^* u \qacc v \tx{ mit } u \in \Gamma^*, v \in \Gamma^+}\]
Wir sagen, dass $M$ eine Funktion $f: \Sigs \ra \Gamma^*$ berechnet, falls f�r alle $x \in \Def{(f)}$
\[q_0 \cent x \vd^* \cent w \qacc f(x)\]
mit $w \in \Gamma^*$ beliebig.
\Mark{Definition 5.5}
\end{fdefinition}
\begin{fsatz}[Church-Turing-These]
Turingmaschinen  sind die Formalisierung des intuitiven Begriffes \ggq{Algorithmus}.
\[\tx{Funktion \textbf{intuitiv} berechenbar} \Lra \tx{Funktion Turing-berechenbar}\]
Indikatoren f�r Korrektheit dieser These
\begin{itemize}
\item Erfahrung aus mehr als 70 Jahren hat gezeigt, dass es (bis jetzt) keine intuitiv berechenbare Funktion gibt, die nicht Turing-berechenbar ist.
\item Turings Analyse des Rechenvorgangs scheint umfassend
\item andere Zug�nge zum Berechenbarkeitsbegriff ($\lambda$-Kalk�l, $\mu$-rekursive Funktionen) sind �quivalent
\item alle modernen Programmiersprachen lassen sich auf \ac{TM} reduzieren
\end{itemize}
\Mark{Satz 5.1}
\end{fsatz}
\begin{fdefinition}
Eine Sprache $L$ hei�t rekursiv aufz�hlbar, wenn eine \ac{TM} $M$ mit $L(M) = L$ existiert.
\[\mathcal{L}_{\tx{RE}} = \gklamm{L(M) \vert M \tx{ ist \ac{TM}}}\]
ist die Menge aller rekursiv aufz�hlbarer Sprachen.
\Mark{Definition 5.6}
\end{fdefinition}
\begin{fdefinition}
Eine Sprache $L$ hei�t \textbf{rekursiv} (entscheidbar), falls eine \ac{TM} $M$ existiert, so dass $\forall w \in \Sigs$ gilt:
\begin{itemize}
\item $q_0 \cent w \vd^* u \qacc v$, falls $w \in L$
\item $q_0 \cent w \vd^* u \qrej v$, fa,,s $w \notin L$
\end{itemize}
mit $u \in \Gamma^*$, $v \in \Gamma^+$.
\[\mathcal{L}_{\tx{R}} = \gklamm{L (M) \vert M \tx{ ist \ac{TM}, die immer h�lt}}\]
ist die Menge der rekursiven \ac{bzw.} algorithmisch erkennbaren Sprachen.\\
Fragen:
\begin{itemize}
\item Gibt es Sprachen, die nicht in $L_{\tx{RE}}$ liegen?
\item Ist $L_{\tx{R}} = L_{\tx{RE}}$?
\end{itemize}
\begin{center}
\begin{tikzpicture}
\draw (0,0) circle (0.5cm);
\draw (0,0) circle (1cm);
\draw (0,0) circle (1.5cm);
\node at (0,0) {$L_{\tx{R}}$};
\node at (0,-1cm) [above] {$L_{\tx{RE}}$};
\end{tikzpicture}
\end{center}
\Mark{Definition 5.7}
\end{fdefinition}

\section{Berechenbarkeit}
\Mark{Section 5.5}
\subsection{Die Methode der Diagonalisierung}
\Mark{Subsection 5.5.1}
\begin{ziel}
Nachweis der Existenz von Sprachen, die nicht rekursiv aufz�hlbar sind.
\end{ziel}
Um dies zu zeigen werden wir folgende quantitative �quivalente Aussage beweisen.
\begin{quote}
\ggq{Die M�chtigkeit der Menge aller \ac{TM} ist kleiner als die M�chtigkeit aller Sprachen �ber $\SigB$.}
\end{quote}
Kodierung von Turingmaschinen (Ziel: Nummerierung \ac{bzw.} Ordnung von Turingmaschinen).\\
Sei $M = \rkl{Q, \Sigma, \Gamma, \delta, q_0, \qacc, \qrej}$ eine \ac{TM} mit
\begin{itemize}
\item $Q = \gklamm{q_1, \dots, q_m, q_0, \qacc, \qrej}$
\item $\Gamma = \gklamm{a_1, \dots, a_r}$
\end{itemize}
Die Kodierung der einzelnen Symbole legen wir wie folgt fest
\begin{multicols}{2}
\begin{itemize}
\item $\Code{(q_i)} = 1 0^i 1$ f�r $i = 1 \dots m$
\item $\Code{(q_0)} = 1 0^{m + 1} 1$
\item $\Code{(\qacc)} = 1 0^{m + 2} 1$
\item $\Code{(\qrej)} = 1 0^{m + 3} 1$
\item $\Code{(a_j)} = 11 0^j 11$ $j = 1 \dots r$
\item $\Code{(N)} = 111 0 111$
\item $\Code{(R)} = 111 00 111$
\item $\Code{(L)} = 111 000 111$
\end{itemize}
\end{multicols}
Transitionen werden wie folgt kodiert
\begin{itemize}
\item $\Code{\rkl{\delta (p, a) = (q, b, \alpha)}} =$
		\[\sharp \Code{(p)} \Code{(a)} \Code{(q)} \Code{(b)} \Code{(\alpha)} \sharp\]
\end{itemize}
Eine vollst�ndige \ac{TM} kann dann so kodiert werden:
\[\Code{(M)} = \sharp 0^{m + 3} \sharp 0^r \sharp \Code{(\tx{Transition 1})} \sharp \Code{(\tx{Transition 2})} \dots\]
\begin{fdefinition}
F�r jede \ac{TM} $M$ bezeichnet
\[\Kod{(M)} = \mathcal{h} \rkl{\Code{(M)}}\]
die Kodierung von $M$, wobei $\mathcal{h}$ Homomorphismus mit $\mathcal{h} (0) = 00$, $\mathcal{h} (1) = 11$ und $\mathcal{\sharp} = 10$.
\[\Kod{(TM)} = \gklamm{\Kod{(M)} \vert M \tx{ ist \ac{TM}}}\]
bezeichnet die Menge der Kodierungen aller \ac{TM}.
\end{fdefinition}

%\chapter{�bungen und Tutorium}
\section{�bungen}
\subsection{Endliche Automaten}
\begin{description}
\item[Ziel (fern)] Minimales und vollst�ndiges Modell f�r Computer \ac{bzw.} \textbf{Algorithmus}
\item[Hier]
		\begin{itemize}
		\item Modell f�r \textbf{spezielle} Entscheidungsprobleme
		\item Einf�hrung wichtiger Begriffe
		\end{itemize}
\end{description}

\Insert{GI(U)-16.04.2009-IMG-1}

\begin{tikzpicture}[->,>=stealth',shorten >=1pt,auto,node distance=2.8cm,semithick]
\node[state,initial,initial text=,accepting] (n1) {even $0$};
\node[state] (n2) [right of=n1] {odd $0$};

\path[->](n1) edge[loop above] node[above] {$1$} ();
\path[->](n2) edge[loop above] node[above] {$1$} ();

\path[->](n1) edge[bend left] node[above] {$0$} (n2);
\path[->](n2) edge[bend left] node[below] {$0$} (n1);
\end{tikzpicture}
\Img{GI(U)-16.04.2009-IMG-2}

$L_1 = \gklamm{w \in \Sigma_{\tx{Bool}}^* \vert \betrag{w}_0 \tx{ gerade } \und \betrag{w}_1 \tx{ gerade}}$

\Insert{GI(U)-16.04.2009-IMG-3}

$L_2 = \gklamm{w \in \Sigma_{\tx{Bool}}^* \vert \betrag{w}_0 \tx{ gerade } \und \betrag{w}_1 \tx{ gerade} \oder \betrag{w}_0 \tx{ gerade } \und \betrag{w}_1 \tx{ ungerade}}$

\Insert{GI(U)-16.04.2009-IMG-4}

$L_3 = \gklamm{w \in \Sigma_{\tx{Bool}}^* \vert Zahlenwert (w) \mod 4 = 0}$
$\approx \gklamm{w \in \Sigma_{\tx{Bool}}^* \vert w \tx{ endet auf } 00}$
Erstmal: $w$ beginnt mit $00$

\Insert{GI(U)-16.04.2009-IMG-5}

\Insert{GI(U)-16.04.2009-IMG-6}

$L_4 = \gklamm{w \in \Sigma_{\tx{Bool}}^* \vert w \tx{ enth�lt } 110110}$

\Insert{GI(U)-16.04.2009-IMG-7}

\section{Tutorium}
\subsection{Aufgaben}
\subsubsection{Vom 21.04.2009}
\begin{enumerate}
\item $L_1 \gklamm{w \in \Sigma_{\tx{Bool}}^* \vert \betrag{w}_1 \mod 2 = 0 \und \betrag{w}_0 = 1}$
\item $L_2 \gklamm{w \in \Sigma_{\tx{Bool}}^* \vert w \tx{ enth�lt } 101101}$
\item $L_3 \gklamm{w \in \Sigma_{\tx{Bool}}^* \vert w \tx{ enth�lt mindestens zweimal } 0}$
\item $L_4 \gklamm{w \in \Sigma_{\tx{Bool}}^* \vert w 0 v \und v \tx{ enth�lt gerade Anzahl an 1en}}$
\end{enumerate}

\subsection{L�sungen}
\subsubsection{Vom 21.04.2009}
\begin{enumerate}
\item L�sung siehe Abbildung \vref{fig:Tutorium_22_01_2009_1}.
		\begin{figure}[htb]
		\centering
		\begin{tikzpicture}[->,>=stealth',shorten >=1pt,auto,node distance=2.8cm,semithick]
		\node[state]
		(node5) 									{DS};
		\node[state,accepting]
		(node3) [above left of=node5] 	{};
		\node[state]
		(node4) [above right of=node5] 	{};
		\node[state,initial,accepting,initial text=]
		(node1) [above of=node3]			{};
		\node[state]
		(node2) [above of=node4] 			{};
		
		\path[->] (node1) edge [bend left] node[above] {$1$} (node2);
		\path[->] (node2) edge [bend left] node[below] {$1$} (node1);
		\path[->] (node1) edge node[right] {$0$} (node3);
		\path[->] (node2) edge node[right] {$0$} (node4);
		\path[->] (node3) edge [bend left] node[above] {$1$} (node4);
		\path[->] (node4) edge [bend left] node[below] {$1$} (node3);
		\path[->] (node3) edge node[left] {$0$} (node5);
		\path[->] (node4) edge node[right] {$0$} (node5);
		\path[->] (node5) edge [loop below] node[right] {$0,1$} ();
		\end{tikzpicture}
		\caption{L�sung zu $L_1$}
		\label{fig:Tutorium_22_01_2009_1}
		\end{figure}

\item L�sung siehe Abbildung \vref{fig:Tutorium_22_01_2009_2}.
		\begin{figure}[htb]
		\centering
		\begin{tikzpicture}[->,>=stealth',shorten >=1pt,auto,node distance=2.8cm,semithick]
		\node[state,initial, initial text=] (n1) {$\epsilon$};
		\node[state] (n2) [right of=n1] {$1$};
		\node[state] (n3) [right of=n2] {$10$};
		\node[state] (n4) [right of=n3] {$101$};
		\node[state] (n5) [right of=n4] {$1011$};
		\node[state] (n6) [below of=n5] {$10110$};
		\node[state,accepting] (n7) [below of=n6] {$101101$};
		
		\path[->] (n1) edge node[above] {$1$} (n2);
		\path[->] (n2) edge node[above] {$0$} (n3);
		\path[->] (n3) edge node[above] {$1$} (n4);
		\path[->] (n4) edge node[above] {$1$} (n5);
		\path[->] (n5) edge node[right] {$0$} (n6);
		\path[->] (n6) edge node[right] {$1$} (n7);
		
		\path[->] (n1) edge [loop below] node[right] {$0$} ();
		\path[->] (n2) edge [loop below] node[right] {$1$} ();
		\path[->] (n3) edge [bend right] node[above] {$0$} (n1);
		\path[->] (n4) edge [bend left] node[below] {$0$} (n3);
		\path[->] (n5) edge [bend left] node[below] {$1$} (n2);
		\path[->] (n6) edge [bend left] node[below] {$1$} (n1);
		\path[->] (n7) edge [loop below] node[right] {$0,1$} ();
		\end{tikzpicture}
		\caption{L�sung zu $L_2$}
		\label{fig:Tutorium_22_01_2009_2}
		\end{figure}

\item L�sung siehe Abbildung \vref{fig:Tutorium_22_01_2009_3}.
		\begin{figure}[htb]
		\centering
		\begin{tikzpicture}[->,>=stealth',shorten >=1pt,auto,node distance=2.8cm,semithick]
		\node[state,initial,initial text=] (n1) {};
		\node[state] (n2) [right of=n1] {};
		\node[state,accepting] (n3) [right of=n2] {};
		
		\path[->] (n1) edge node[above] {$0$} (n2);
		\path[->] (n2) edge node[above] {$0$} (n3);
		
		\path[->] (n1) edge [loop above] node[right] {$1$} ();
		\path[->] (n2) edge [loop above] node[right] {$1$} ();
		\path[->] (n3) edge [loop above] node[right] {$0,1$} ();
		\end{tikzpicture}
		\caption{L�sung zu $L_3$}
		\label{fig:Tutorium_22_01_2009_3}
\end{figure}
\end{enumerate}

%</Text>%%%%%%%%%%%%%%%%%%%%%%%%%%%%%%%%%%%%%%%%%%%%%%%%%%%%%%%%%%%%%%%%%%%%%%%
\backmatter
%Literaturverzeichnis
%\newpage
%\addcontentsline{toc}{part}{Literaturverzeichnis}
%\bibliography{literatur}{}

%Stichwortverzeichnis
\newpage
\renewcommand{\indexname}{Stichwortverzeichnis}
\addcontentsline{toc}{part}{Stichwortverzeichnis}
\printindex

\end{document}
